% Documento de Requisitos do IntegraCASD, versão 1.2.
% Compilar a partir desta pasta (docs/requisitos), duas vezes para o sumário:
%   pdflatex IntegraCASD_DocRequisitos_v1.2.tex
% As figuras vêm de docs/telas/prints/ (protótipo navegável em docs/telas/).

\documentclass[11pt,a4paper,oneside]{report}

\usepackage[utf8]{inputenc}
\usepackage[T1]{fontenc}
\usepackage[brazilian]{babel}
\usepackage{lmodern}
\usepackage[a4paper,top=2.5cm,bottom=2.5cm,left=2.5cm,right=2.5cm,headheight=14pt]{geometry}
\usepackage{amssymb}
\usepackage{graphicx}
\usepackage{xcolor}
\usepackage{array}
\usepackage{tabularx}
\usepackage{longtable}
\usepackage{booktabs}
\usepackage{enumitem}
\usepackage{fancyhdr}
\usepackage{titlesec}
\usepackage[font=small,labelfont=bf]{caption}
\usepackage[section]{placeins}
\usepackage[hidelinks]{hyperref}

\graphicspath{{../telas/prints/}}

% ---------------------------------------------------------------------------
% Aparência
% ---------------------------------------------------------------------------

\definecolor{casdamarelo}{HTML}{FFC20E}
\definecolor{casdpreto}{HTML}{0E0E0E}
\definecolor{cinza}{HTML}{5A5A5A}

\setlength{\parindent}{0pt}
\setlength{\emergencystretch}{3em}
\setlength{\parskip}{0.6em}
\setlist{itemsep=0.2em,topsep=0.3em}

\titleformat{\chapter}[display]
  {\normalfont\huge\bfseries\color{casdpreto}}
  {\filleft\large\color{cinza}\chaptertitlename\ \thechapter}
  {0.5em}
  {\filleft}
  [\vspace{0.3em}{\color{casdamarelo}\titlerule[2pt]}]
\titleformat{name=\chapter,numberless}[display]
  {\normalfont\huge\bfseries\color{casdpreto}}
  {}
  {0pt}
  {\filleft}
  [\vspace{0.3em}{\color{casdamarelo}\titlerule[2pt]}]
\titlespacing*{\chapter}{0pt}{0pt}{2em}
\titleformat{\section}{\normalfont\Large\bfseries\color{casdpreto}}{}{0pt}{}
\titleformat{\subsection}{\normalfont\large\bfseries}{}{0pt}{}

\pagestyle{fancy}
\fancyhf{}
\renewcommand{\chaptermark}[1]{\markboth{#1}{}}
\fancyhead[L]{\small Documento de Requisitos -- IntegraCASD}
\fancyhead[R]{\small\leftmark}
\fancyfoot[L]{\small Versão 1.2}
\fancyfoot[C]{\small\thepage}
\fancyfoot[R]{\small Novembro/2026}
\fancypagestyle{plain}{%
  \fancyhf{}
  \fancyhead[L]{\small Documento de Requisitos -- IntegraCASD}
  \fancyfoot[L]{\small Versão 1.2}
  \fancyfoot[C]{\small\thepage}
  \fancyfoot[R]{\small Novembro/2026}}

% ---------------------------------------------------------------------------
% Macros dos requisitos
% ---------------------------------------------------------------------------

% Cabeçalho de um requisito: \req{RF001}{Manter cadastro de alunos}
\newcommand{\req}[2]{%
  \subsection*{[#1] #2}%
  \addcontentsline{toc}{subsection}{[#1] #2}}

% Requisito retirado do escopo: mantém o identificador reservado.
\newcommand{\reqretirado}[3]{%
  \subsection*{[#1] #2 \textmd{\textit{(retirado do MVP)}}}%
  \addcontentsline{toc}{subsection}{[#1] #2 (retirado do MVP)}%
  \begin{quote}\small\color{cinza}#3\end{quote}}

\newcommand{\marcado}{$\boxtimes$}
\newcommand{\vazio}{$\square$}
\newcommand{\prioE}{\par\textbf{Prioridade:}\quad\marcado~Essencial\quad\vazio~Importante\quad\vazio~Desejável\par}
\newcommand{\prioI}{\par\textbf{Prioridade:}\quad\vazio~Essencial\quad\marcado~Importante\quad\vazio~Desejável\par}
\newcommand{\prioD}{\par\textbf{Prioridade:}\quad\vazio~Essencial\quad\vazio~Importante\quad\marcado~Desejável\par}

\newcommand{\ator}[1]{\par\textbf{Ator:} #1\par}
\newcommand{\entradas}{\textbf{Entradas e pré-condições:} }
\newcommand{\saidas}{\par\textbf{Saídas e pós-condições:} }
\newcommand{\fluxoprincipal}{\par\textbf{Fluxo de eventos principal}\par}
\newcommand{\fluxossecundarios}{\par\textbf{Fluxos secundários (alternativos e de exceção)}\par}
\newcommand{\fluxo}[2]{\par\smallskip\textit{#1 -- #2}\par}
\newcommand{\casos}[1]{\par{\raggedright\textbf{Caso(s) de uso associado(s):} #1\par}}

% Figura de tela: \tela[altura relativa]{arquivo}{legenda}{rótulo}
\newcommand{\tela}[4][0.42]{%
  \begin{figure}[htbp]
    \centering
    \fbox{\includegraphics[width=0.97\linewidth,height=#1\textheight,keepaspectratio]{#2}}
    \caption{#3}
    \label{fig:#4}
  \end{figure}}

\begin{document}

% ---------------------------------------------------------------------------
% Ficha técnica
% ---------------------------------------------------------------------------

\begin{titlepage}
  \pagestyle{empty}
  {\large Departamento de Ciência da Computação, ITA-SJC\par}
  \vspace{0.5em}
  {\color{casdamarelo}\rule{\linewidth}{2pt}\par}
  \vspace{2em}
  {\Huge\bfseries Documento de Requisitos\par}
  \vspace{0.5em}
  {\LARGE IntegraCASD\par}
  \vspace{0.3em}
  {\large Sistema de gestão financeira do Centro Acadêmico Santos Dumont\par}

  \vspace{3em}
  {\large\bfseries Ficha Técnica\par}
  \vspace{1em}
  \textbf{Equipe Responsável pela Elaboração}\par
  \begin{tabular}{@{}ll@{}}
    João Gabriel de Magalhães Gonzaga & ITA-SJC \\
    Gabriel Negreiro Oliveira         & ITA-SJC \\
    Marcello Ryaj de Almeida Santos   & ITA-SJC \\
    Mateus Levi Morais Feitosa        & ITA-SJC \\
  \end{tabular}

  \vspace{1.5em}
  \textbf{Público Alvo}\par
  Este documento destina-se aos integrantes do projeto IntegraCASD (analistas, projetistas,
  programadores e testadores), à equipe financeira do CASD e aos docentes responsáveis pela
  avaliação deste trabalho na disciplina CSI-28 -- Fundamentos de Engenharia de Software.

  \vfill
  \begin{center}
    {\large\bfseries Versão 1.2 -- São José dos Campos, Novembro de 2026\par}
  \end{center}
  \vspace{1em}
  Dúvidas, críticas e sugestões devem ser encaminhadas por escrito para um dos seguintes endereços
  eletrônicos: \texttt{marcelloryaj37@gmail.com} ou \texttt{marcello.santos.9107@ga.ita.br}.
  Recomendamos que o assunto seja identificado com o título desta obra e que a mensagem traga o
  nome completo do remetente, para que seja possível o envio de respostas.
\end{titlepage}

% ---------------------------------------------------------------------------
% Histórico de revisões
% ---------------------------------------------------------------------------

\chapter*{Histórico de revisões}
\markboth{Histórico de revisões}{}

\begin{tabularx}{\linewidth}{@{}l l X@{}}
  \toprule
  \textbf{Versão} & \textbf{Data} & \textbf{Descrição} \\
  \midrule
  1.0 & Setembro/2026 & Versão inicial, com o levantamento completo dos requisitos. \\
  1.1 & Outubro/2026  & Ajuste ao escopo do MVP definido pela equipe: retirada dos requisitos
                        fora do escopo, atualização dos casos de uso com as decisões tomadas desde
                        a versão 1.0 e inclusão das telas do protótipo navegável. \\
  1.2 & Novembro/2026 & Retirada do MVP do reenvio do boleto por e-mail ([Cobranças.RF003]). \\
  \bottomrule
\end{tabularx}

\section*{Alterações em relação à versão 1.1}

\textbf{Requisitos alterados.}
\begin{itemize}
  \item {[Cobranças.RF003]}: o reenvio do boleto por e-mail sai do MVP, e o caso de uso passa a
        cobrir apenas a obtenção do boleto para download. A API do Asaas não oferece operação para
        reenviar a notificação de uma cobrança, e o envio pelo próprio IntegraCASD depende de um
        domínio próprio verificado no provedor de correio eletrônico, que o CASD ainda não possui.
        O reenvio por e-mail fica para depois do MVP. Saem o passo de escolha entre baixar e
        reenviar e o fluxo A2 (reenvio a partir da inadimplência).
  \item {[Indicadores.RF002]}: no fluxo A1, a ação a partir da relação de devedores passa a ser
        abrir as cobranças vencidas do aluno e obter o boleto, em vez de reenviá-lo.
\end{itemize}

\textbf{Telas.} As figuras do detalhe da cobrança, do painel de indicadores e da aba de cobranças
do aluno foram atualizadas para refletir a retirada do reenvio.

\section*{Alterações em relação à versão 1.0}

\textbf{Requisitos retirados do MVP.} Os identificadores abaixo permanecem reservados e não são
reaproveitados, para que as referências já feitas no código, no protótipo e no quadro de tarefas
continuem válidas. Cada um aparece no Capítulo~2 apenas com a indicação da retirada.
\begin{itemize}
  \item {[Cadastros.RF003]} Manter cadastro de fornecedores e parceiros. O contrato passa a guardar
        apenas o nome da contraparte, em texto.
  \item {[Movimentações.RF004]} Anexar comprovante a uma movimentação.
  \item {[Indicadores.RF004]} Exportar relatório.
\end{itemize}

\textbf{Outros itens retirados do escopo.} O cadastro de assinaturas internas (por exemplo, a
lavanderia) deixa de fazer parte do MVP; os valores correspondentes são lançados como
movimentações de entrada classificadas por categoria.

\textbf{Requisitos alterados.}
\begin{itemize}
  \item {[Cadastros.RF001]}: as entradas do cadastro de alunos seguem o modelo de dados definido pela
        equipe (turma, curso, posto, matrícula, alojamento, celular e apelido); a data de
        nascimento deixa de ser coletada, por não ser exigida pelo Asaas
        ([Padrões.NF001]). Foi incluído o fluxo de cadastro incompleto vindo da carga inicial.
  \item {[Cadastros.RF002]}: inclusão do salário de referência do funcionário.
  \item {[Cadastros.RF004]}: o contrato deixa de depender do cadastro de fornecedores e passa a
        apresentar o total pago no ano.
  \item {[Cobranças.RF006]}: o caso de uso passa a se chamar ``Registrar pagamento manualmente'' e
        cobre também o pagamento combinado diretamente com a área administrativa.
  \item {[Cobranças.RF007]}: inclusão da sincronização de uma única cobrança e da execução periódica
        automática.
  \item {[Movimentações.RF001]}: o vínculo com fornecedor é substituído pelo vínculo com contrato;
        inclusão da forma de pagamento ``Transferência''.
  \item {[Administração.RF003]} e [Administração.RF004]: inclusão do histórico de execuções da
        geração recorrente e da configuração do prazo de alerta de interrupção da conciliação.
  \item {[Desempenho.NF002]}, [Padrões.NF001], [Padrões.NF003], [Hardware e software.NF001] e
        [Hardware e software.NF005]: atualização conforme as decisões de arquitetura já
        registradas.
\end{itemize}

\textbf{Requisitos incluídos.} [Usabilidade.NF004] Informação que não depende apenas de cor.

\textbf{Capítulo incluído.} Capítulo~4, com a descrição da interface com o usuário. As telas do
protótipo também aparecem junto aos casos de uso a que se referem, no Capítulo~2.

\tableofcontents

% ---------------------------------------------------------------------------
% Introdução
% ---------------------------------------------------------------------------

\chapter*{Introdução}
\addcontentsline{toc}{chapter}{Introdução}
\markboth{Introdução}{}

Este documento especifica o sistema IntegraCASD, fornecendo aos desenvolvedores as informações
necessárias para o projeto e a implementação, assim como para a realização dos testes e a
homologação do sistema.

O IntegraCASD destina-se à gestão financeira e administrativa do Centro Acadêmico Santos Dumont
(CASD). Atualmente, essas informações encontram-se distribuídas em diferentes planilhas, mantidas
de forma independente, o que dificulta o cruzamento de dados, torna a consulta lenta e sujeita a
erros e impede uma visão consolidada da situação financeira do Centro Acadêmico. O sistema aqui
especificado tem por finalidade centralizar esses dados em um único ambiente, automatizar a emissão
e a conciliação de cobranças e disponibilizar indicadores que apoiem a tomada de decisão pela
administração e pela diretoria do CASD.

Desde a versão 1.1, a especificação se restringe ao escopo do produto mínimo viável (MVP), a ser
entregue em dois meses. Os requisitos que ficaram de fora estão indicados no Histórico de revisões.

\section*{Visão geral deste documento}
\addcontentsline{toc}{section}{Visão geral deste documento}

Esta introdução fornece as informações necessárias para fazer um bom uso deste documento,
explicitando seus objetivos e as convenções que foram adotadas no texto, além de conter uma lista
de referências para outros documentos relacionados. As demais partes apresentam a especificação do
sistema IntegraCASD e estão organizadas como descrito abaixo.
\begin{itemize}
  \item \textbf{Capítulo 1 -- Descrição geral do sistema:} apresenta uma visão geral do sistema,
        caracterizando qual é o seu escopo e descrevendo seus usuários.
  \item \textbf{Capítulo 2 -- Requisitos funcionais (casos de uso):} especifica os requisitos
        funcionais do sistema, descrevendo os fluxos de eventos, prioridades, atores, entradas e
        saídas de cada caso de uso, acompanhados das telas do protótipo que os ilustram.
  \item \textbf{Capítulo 3 -- Requisitos não funcionais:} especifica os requisitos não funcionais
        do sistema, divididos em requisitos de usabilidade, confiabilidade, desempenho, segurança,
        distribuição, adequação a padrões e requisitos de hardware e software.
  \item \textbf{Capítulo 4 -- Descrição da interface com o usuário:} apresenta a organização das
        telas, as decisões de interface que valem para todas elas e a correspondência entre telas e
        casos de uso.
  \item \textbf{Apêndice A -- Pontos em aberto:} lista as questões que ainda dependem de decisão
        da equipe ou do CASD.
\end{itemize}

\section*{Convenções, termos e abreviações}
\addcontentsline{toc}{section}{Convenções, termos e abreviações}

A correta interpretação deste documento exige o conhecimento de algumas convenções e termos
específicos, que são descritos a seguir.

\begin{longtable}{@{}>{\bfseries\raggedright\arraybackslash}p{0.24\linewidth} p{0.70\linewidth}@{}}
  \toprule
  \textmd{\textbf{Termo}} & \textbf{Definição} \\
  \midrule
  \endhead
  Asaas & Provedor externo de serviços financeiros contratado pelo CASD, responsável pela emissão
          das cobranças e pela comunicação com a instituição bancária. \\
  Baixa & Registro, no sistema, de que uma cobrança foi quitada, com a consequente atualização da
          situação do pagador. \\
  Boleto & Documento de cobrança emitido em favor do CASD, identificado por um número único, com
           valor e data de vencimento definidos. \\
  CASD & Centro Acadêmico Santos Dumont. \\
  Categoria & Classificação atribuída a uma movimentação financeira, utilizada para agrupar
              entradas e saídas na apuração dos indicadores. \\
  Cobrança & Solicitação de pagamento registrada no sistema, dirigida a um aluno e sempre
             correspondente a uma cobrança criada no Asaas. \\
  Conciliação & Atualização automática da situação de uma cobrança a partir das notificações do
                Asaas, com a criação da movimentação de entrada correspondente ao pagamento. \\
  Contraparte & Pessoa ou empresa com a qual o CASD firma um contrato. No MVP é registrada apenas
                pelo nome, sem cadastro próprio. \\
  Estorno & Lançamento de sinal contrário que anula o efeito de uma movimentação, preservando o
            registro original. \\
  Inadimplente & Aluno que possui ao menos uma cobrança cuja data de vencimento já transcorreu sem
                 que o pagamento tenha sido identificado. \\
  Movimentação & Todo registro de entrada ou de saída de recursos financeiros do CASD, incluindo
                 recebimentos de alunos, pagamentos de contratos, despesas de pessoal e gastos com
                 projetos. \\
  MVP & Produto mínimo viável: o conjunto de funcionalidades da primeira versão entregue. \\
  Pagador & Pessoa a quem uma cobrança é dirigida, tipicamente um aluno do CASD. \\
  Situação da cobrança & Uma das seguintes: Aberta, Vencida, Paga, Paga com registro manual,
                         Cancelada ou Estornada. \\
  Sincronização & Consulta ao Asaas que atualiza as cobranças cujas notificações não chegaram ao
                  sistema. \\
  Webhook & Notificação enviada automaticamente pelo Asaas ao IntegraCASD sempre que ocorre uma
            mudança na situação de uma cobrança, empregada pelo sistema para realizar a
            conciliação. \\
  \bottomrule
\end{longtable}

\subsection*{1. Identificação dos requisitos}

Por convenção, a referência a requisitos é feita através do nome da subseção onde eles estão
descritos, seguido do identificador do requisito, de acordo com o esquema abaixo:
\begin{center}
  [nome da subseção.identificador do requisito]
\end{center}
Por exemplo, o requisito [Cobranças.RF005] está descrito na subseção ``Cobranças'' do
Capítulo~2, em um bloco identificado por [RF005]. Já o requisito não funcional
[Confiabilidade.NF004] está descrito na seção de requisitos não funcionais de Confiabilidade, em um
bloco identificado por [NF004]. Os identificadores de requisitos retirados do escopo não são
reaproveitados.

\subsection*{2. Prioridades dos requisitos}

Para estabelecer a prioridade dos requisitos foram adotadas as denominações ``essencial'',
``importante'' e ``desejável''.
\begin{itemize}
  \item \textbf{Essencial} é o requisito sem o qual o sistema não entra em funcionamento.
        Requisitos essenciais são imprescindíveis e têm que ser implementados impreterivelmente.
  \item \textbf{Importante} é o requisito sem o qual o sistema entra em funcionamento, mas de forma
        não satisfatória. Requisitos importantes devem ser implementados, mas, se não forem, o
        sistema poderá ser implantado e usado mesmo assim.
  \item \textbf{Desejável} é o requisito que não compromete as funcionalidades básicas do sistema,
        isto é, o sistema pode funcionar de forma satisfatória sem ele. Requisitos desejáveis
        podem ser deixados para versões posteriores, caso não haja tempo hábil para implementá-los
        na versão que está sendo especificada.
\end{itemize}

\section*{Referências}
\addcontentsline{toc}{section}{Referências}

\begin{enumerate}[before=\raggedright]
  \item Lei nº 13.709, de 14 de agosto de 2018 (Lei Geral de Proteção de Dados Pessoais);
        Presidência da República, Casa Civil;
        \url{https://www.planalto.gov.br/ccivil_03/_ato2015-2018/2018/lei/l13709.htm}
  \item Documentação da interface de programação do Asaas, seções de Cobranças e de Webhooks;
        consultada em setembro de 2026; Asaas Gestão Financeira S.A.;
        \url{https://docs.asaas.com/}
  \item IntegraCASD -- Documento de Requisitos, versão 1.0; setembro de 2026;
        \url{docs/requisitos/IntegraCASD_DocRequisitos_v1.pdf}.
  \item IntegraCASD -- Relatório Descritivo do Projeto, versão 1.0 (Laboratório 03); setembro de
        2026; \url{docs/requisitos/IntegraCASD_RelDescritivo_v1.pdf}.
  \item Registros de decisões de arquitetura: ADR-0001, Monólito modular como estilo arquitetural;
        ADR-0002, PostgreSQL como banco e como fila; outubro de 2026; \texttt{docs/adr/}.
  \item Diagramas de arquitetura (C4 e implantação) e diagrama de sequência da conciliação;
        \texttt{docs/arquitetura/}.
  \item Protótipo navegável das telas, com dados fictícios; outubro de 2026;
        \texttt{docs/telas/}.
\end{enumerate}

% ---------------------------------------------------------------------------
% Capítulo 1
% ---------------------------------------------------------------------------

\chapter{Descrição geral do sistema}

O IntegraCASD tem por objetivo centralizar, em um único ambiente, a gestão financeira e
administrativa do Centro Acadêmico Santos Dumont, hoje distribuída entre planilhas mantidas de
forma independente. O sistema reúne o cadastro dos alunos e dos funcionários, o registro dos
contratos e dos projetos, o registro de todas as movimentações de entrada e saída de recursos, a
emissão e a conciliação das cobranças dirigidas aos alunos e a apuração dos indicadores que
descrevem a situação financeira do Centro Acadêmico.

O propósito da aplicação não é apenas substituir as planilhas por um banco de dados. O ganho
pretendido está no relacionamento entre informações que hoje vivem separadas: saber que determinado
aluno possui mensalidades em atraso sem conferir o extrato bancário linha a linha; saber quanto foi
gasto em um projeto ao longo de um semestre sem precisar somar linhas de três arquivos diferentes;
saber quantos alunos estão inadimplentes sem depender de conferência manual. É esse cruzamento que o
processo atual não oferece e que motiva o desenvolvimento do sistema.

O projeto é relevante para todos os envolvidos. Para a responsável pela área administrativa,
porque elimina o trabalho repetitivo de emitir cobranças uma a uma e de conferir manualmente quais
foram pagas. Para a diretoria do CASD, porque disponibiliza indicadores confiáveis e atualizados
para decidir sobre gastos, contratos e projetos. Para o próprio Centro Acadêmico como instituição,
porque preserva a memória financeira entre gestões, hoje dependente da transferência informal de
arquivos a cada troca de diretoria.

Em linhas gerais, o IntegraCASD é uma aplicação web de três camadas: uma interface de usuário
acessível por navegador, um servidor de aplicação que concentra as regras de negócio e um banco de
dados relacional que armazena o cadastro, as movimentações e as cobranças. O servidor é organizado
como um monólito modular, com um módulo para cada subseção do Capítulo~2, conforme o ADR-0001.

\section{Abrangência e sistemas relacionados}

O MVP do IntegraCASD contempla o cadastro de alunos e de funcionários; o registro de contratos e de
projetos; o registro das movimentações financeiras de entrada e de saída, classificadas por
categoria; a geração automática das cobranças recorrentes e a emissão manual de cobranças avulsas;
a emissão do boleto e o seu envio ao pagador; a conciliação automática dos pagamentos; e a
apresentação de um painel de indicadores com a posição financeira do Centro Acadêmico, o total de
inadimplentes, a relação dos maiores devedores e a distribuição dos gastos por categoria. Essas
funcionalidades são apenas citadas aqui, para dar ao leitor uma ideia geral dos serviços oferecidos;
os detalhes constam do Capítulo~2.

Estão fora do MVP, podendo ser retomados em versões futuras:
\begin{itemize}
  \item \textbf{Cadastro de fornecedores e de organizações parceiras.} O contrato registra apenas o
        nome da contraparte, em texto ([Cadastros.RF004]).
  \item \textbf{Assinaturas internas.} Serviços como a lavanderia não têm cadastro próprio nem
        geram cobrança; seus valores são lançados como movimentações classificadas por categoria.
  \item \textbf{Anexo de comprovante às movimentações.}
  \item \textbf{Exportação de relatórios} em PDF ou CSV.
\end{itemize}

Estão, ainda, explicitamente fora do escopo do sistema:
\begin{itemize}
  \item \textbf{Acesso de alunos e de membros do CASD ao sistema.} O aluno é tratado como pagador,
        e não como usuário: ele recebe a cobrança por e-mail e efetua o pagamento fora do
        IntegraCASD, sem realizar acesso autenticado. A disponibilização de uma área em que o
        próprio aluno consulte sua situação financeira foi avaliada e adiada para uma versão
        futura.
  \item \textbf{Processamento de pagamentos.} O sistema não recebe, não movimenta e não custodia
        recursos financeiros. A emissão da cobrança e o recebimento do valor são realizados
        integralmente pelo Asaas; ao IntegraCASD cabe solicitar a cobrança e registrar o seu
        desfecho.
  \item \textbf{Folha de pagamento e obrigações trabalhistas.} O sistema registra os pagamentos
        efetuados a funcionários como movimentações de saída, mas não calcula salários, encargos,
        férias ou tributos, atribuições que permanecem a cargo da contabilidade contratada pelo
        CASD.
  \item \textbf{Escrituração contábil e prestação de contas oficial.} O IntegraCASD é uma
        ferramenta de gestão, e não um sistema contábil. Os indicadores que produz servem à
        decisão interna e não substituem os demonstrativos elaborados pela contabilidade.
  \item \textbf{Emissão de notas fiscais}, por não integrar a rotina hoje executada pela área
        administrativa.
  \item \textbf{Controle de estoque e de patrimônio.} Compras de materiais e bens são registradas
        apenas como movimentações financeiras; o acompanhamento físico dos itens não faz parte do
        escopo.
\end{itemize}

O IntegraCASD não é um sistema autocontido: sua operação depende de um sistema externo, descrito a
seguir.

\begin{tabularx}{\linewidth}{@{}l X@{}}
  \toprule
  \textbf{Sistema} & \textbf{Papel e meio de comunicação} \\
  \midrule
  Asaas & Provedor de serviços financeiros responsável por emitir as cobranças em favor do CASD,
          disponibilizar o boleto ao pagador, enviá-lo por e-mail e receber o pagamento. O
          IntegraCASD solicita a criação e o cancelamento das cobranças por meio da API REST do
          Asaas e recebe, por webhook, as notificações de mudança de situação de cada cobrança, que
          utiliza para realizar a conciliação. A comunicação ocorre nos dois sentidos, e a
          indisponibilidade do Asaas impede a emissão de novas cobranças, ainda que não impeça o uso
          das demais funções do sistema. No código, apenas o componente Cliente Asaas comunica-se
          com a plataforma; o restante do sistema depende de uma interface. \\
  \bottomrule
\end{tabularx}

O envio do boleto por e-mail ao aluno é realizado pelo próprio Asaas, acionado pelo IntegraCASD no
momento da criação da cobrança. O sistema não mantém servidor de correio próprio para essa
finalidade, limitando-se a registrar que o envio foi solicitado e a oferecer o download do boleto
para os casos em que a área administrativa opte por entregá-lo por outro meio. O reenvio do boleto
por e-mail a partir do IntegraCASD não faz parte do MVP ([Cobranças.RF003]).

As planilhas hoje utilizadas pelo CASD não constituem um sistema relacionado, e sim a origem dos
dados da carga inicial ([Distribuição.NF002]). Após a implantação, elas deixam de ser mantidas e o
IntegraCASD passa a ser a única fonte da informação financeira do Centro Acadêmico. O cadastro de
alunos é criado a partir dessas planilhas e, a partir de então, mantido dentro do sistema pela área
administrativa.

\section{Descrição dos usuários}

O IntegraCASD possui dois tipos de usuário, ambos internos ao Centro Acadêmico e em número
reduzido: a responsável pela área administrativa e os integrantes da diretoria do CASD. Não há uso
pelo público externo, e o volume de acessos simultâneos é pequeno -- característica que afasta
preocupações de escala, mas não reduz a exigência de correção: o sistema é a única fonte da
informação financeira do Centro Acadêmico, e um dado errado nele é um dado errado em toda parte.

O principal problema que limita a produtividade desses usuários hoje é a fragmentação da
informação. Como cada assunto é controlado em uma planilha própria, qualquer pergunta que atravesse
dois arquivos exige conferência manual: identificar quem está devendo obriga a confrontar a
planilha de cobranças com o extrato bancário, linha a linha; apurar o gasto de um projeto obriga a
somar lançamentos dispersos. Além do tempo consumido, o processo é sujeito a erro e não deixa
registro de quem alterou o quê, o que dificulta tanto a auditoria quanto a transferência de
responsabilidades entre gestões.

Um terceiro grupo de pessoas relaciona-se com o sistema sem ser seu usuário: os alunos, na condição
de pagadores. Eles recebem cobranças originadas no IntegraCASD e seus dados são nele armazenados,
mas não realizam acesso autenticado nem interagem com qualquer interface do sistema. Essa distinção
é relevante porque os alunos não geram requisitos de usabilidade, mas geram requisitos de segurança
e de proteção de dados pessoais, tratados no Capítulo~3.

\subsection*{1. Responsável pela área administrativa}

Funcionária do CASD responsável pela rotina administrativa e financeira do Centro Acadêmico. É a
usuária de uso mais intenso do sistema e a que realiza o maior volume de operações: mantém os
cadastros de alunos e funcionários; lança as movimentações de entrada e de saída; emite cobranças
avulsas; acompanha a situação das cobranças recorrentes; e verifica as pendências de conciliação.

Utiliza o sistema diariamente, a partir de um computador, e conhece a fundo o processo financeiro
do CASD, inclusive as exceções e os casos particulares que hoje só existem na sua memória e nas
observações escritas à margem das planilhas. É, por isso, a principal fonte de informação para o
levantamento de requisitos deste projeto. Por outro lado, não possui formação em informática, e o
sistema não pode pressupor familiaridade com vocabulário técnico: mensagens de erro precisam
explicar o que fazer, e não apenas o que falhou.

Sua expectativa central em relação ao sistema é a redução do trabalho repetitivo. As operações que
executa muitas vezes ao dia, como localizar um aluno, lançar uma movimentação ou emitir uma
cobrança, precisam ser rápidas e alcançáveis em poucos passos, ainda que isso implique telas mais
densas do que seria adequado a um usuário eventual.

\subsection*{2. Diretoria do CASD}

Alunos do ITA que ocupam cargos na diretoria do Centro Acadêmico, com destaque para a diretoria
financeira. Utilizam o sistema com frequência menor e propósito distinto: acompanham os indicadores
do painel para decidir sobre gastos, contratos e projetos, e consultam a situação de inadimplência
para orientar as ações de cobrança.

Concentram também as funções de configuração e administração do sistema: cadastro e permissão dos
usuários, definição das categorias de movimentação, parametrização dos valores e das datas de
vencimento das cobranças recorrentes e guarda das credenciais de integração com o Asaas. Essa
separação é deliberada: as decisões que alteram o comportamento do sistema como um todo cabem à
instância responsável pela gestão do Centro Acadêmico, e não a quem executa a rotina diária.

Duas características desse grupo têm consequência direta sobre o projeto. A primeira é a alta
rotatividade: a diretoria é renovada periodicamente, de modo que o sistema recebe usuários novos
com regularidade, sem tempo hábil para treinamento extenso, e as telas de indicadores e de
configuração precisam ser compreensíveis por si mesmas. A segunda é que se trata de usuários com
boa familiaridade com tecnologia, mas não necessariamente com o processo financeiro do CASD,
exatamente o inverso do perfil da área administrativa. O sistema precisa atender bem aos dois.

% ---------------------------------------------------------------------------
% Capítulo 2
% ---------------------------------------------------------------------------

\chapter{Requisitos funcionais (casos de uso)}

Este capítulo apresenta os requisitos funcionais do MVP do IntegraCASD. Os casos de uso foram
agrupados em cinco subseções de requisitos correlacionados, que correspondem também aos módulos do
servidor, e a numeração dos identificadores é reiniciada em cada uma delas, conforme a convenção
descrita na Introdução. Junto a cada caso de uso aparecem as telas do protótipo que o ilustram; os
dados mostrados nelas são fictícios.

Duas convenções valem para o capítulo inteiro. A primeira é que o perfil de diretoria possui todas
as permissões do perfil da área administrativa, acrescidas das funções de configuração reunidas na
subseção Administração. Quando um caso de uso indica a área administrativa como ator, a diretoria
também pode executá-lo. A segunda é que nem todo caso de uso tem ator humano: a geração de
cobranças recorrentes é disparada pela passagem do tempo, e a conciliação de pagamentos é disparada
pelo Asaas. Esses casos de uso estão especificados aqui, e não entre os requisitos não funcionais,
porque descrevem comportamento do sistema, e não qualidades dele.

Valem, ainda, duas regras de dados para todos os casos de uso: valores monetários são registrados
em centavos, como números inteiros, e toda operação que altera dados é registrada na trilha de
auditoria ([Segurança.NF006]).

% ===========================================================================
\section{Cadastros}

Esta subseção reúne os casos de uso de manutenção das pessoas e dos vínculos que dão origem a
cobranças e a despesas. Todos seguem o padrão ``manter'', que reúne inclusão, alteração, consulta e
inativação em um único caso de uso, com as variações descritas nos fluxos secundários. Nenhum
registro cadastral é excluído fisicamente do sistema: a inativação preserva o histórico financeiro
associado.

Os cadastros ficam reunidos em uma única tela, com uma aba para cada tipo de registro.

\textbf{Ator de todos os casos de uso desta subseção:} área administrativa.

% ---------------------------------------------------------------------------
\req{RF001}{Manter cadastro de alunos}

Permite incluir, alterar, consultar e inativar o cadastro dos alunos do CASD, que são os pagadores
de todas as cobranças emitidas pelo sistema.

\prioE

\entradas o usuário está autenticado. Para inclusão, o CPF informado não pode existir em outro
cadastro ativo. Entradas: nome completo, CPF, e-mail, celular, apelido, matrícula, turma, curso,
posto (civil, aluno militar ou aspirante a oficial), alojamento (bloco, apartamento e vaga) e
situação (ativo ou inativo). O nome, o CPF e o e-mail são obrigatórios, sendo os dois últimos
exigidos pelo Asaas para a emissão da cobrança.

\saidas o cadastro é criado, alterado ou inativado. Aluno ativo passa a ser considerado na geração
de cobranças recorrentes ([Cobranças.RF001]); aluno inativo deixa de sê-lo, mas mantém visíveis as
cobranças e as movimentações já registradas em seu nome.

\fluxoprincipal
\begin{enumerate}
  \item O usuário seleciona a opção ``Cadastros'', aba ``Alunos'', ou usa a busca da barra
        superior, e informa o nome ou o CPF a ser pesquisado (Figura~\ref{fig:cadastros-alunos}).
  \item O sistema apresenta os cadastros correspondentes, com a situação e o valor em aberto de
        cada aluno.
  \item O usuário seleciona um aluno para abrir o seu cadastro, ou aciona ``Novo aluno''.
  \item O usuário preenche ou altera os dados cadastrais (Figuras~\ref{fig:cadastros-novo-aluno}
        e~\ref{fig:aluno-dados}).
  \item O sistema valida os campos obrigatórios e a consistência do CPF e do e-mail.
  \item O sistema grava o cadastro e registra a operação na trilha de auditoria.
\end{enumerate}

\fluxossecundarios

\fluxo{A1}{Inativação de aluno}
O usuário aciona a opção ``Inativar aluno'' e confirma. O sistema altera a situação do aluno, que
deixa de gerar novas cobranças recorrentes, e mantém o cadastro consultável.

\fluxo{A2}{Inclusão a partir da busca}
Se a busca não encontrar nenhum cadastro, o sistema oferece a inclusão de um novo aluno, já com o
nome ou o CPF pesquisado preenchido.

\fluxo{A3}{Consulta do histórico do aluno}
O cadastro do aluno apresenta, em abas, as cobranças em seu nome, com o total em aberto
(Figura~\ref{fig:aluno-cobrancas}), e as movimentações de entrada geradas por elas
(Figura~\ref{fig:aluno-movimentacoes}). A partir do cadastro, o usuário pode emitir uma cobrança
avulsa com o aluno já preenchido ([Cobranças.RF002]). As movimentações exibidas ali não são
corrigidas a partir do cadastro: a correção é feita pela própria cobrança.

\fluxo{A4}{Cadastro incompleto}
Alunos vindos da carga inicial podem não ter e-mail ou CPF válido. A lista sinaliza esses cadastros
(``Falta e-mail''), porque eles não recebem cobrança enquanto não forem corrigidos
([Cobranças.RF001], fluxo A2).

\fluxo{E1}{CPF já cadastrado}
Se o CPF informado já pertencer a outro cadastro ativo, o sistema não cria o novo registro e
apresenta o cadastro existente para eventual atualização, evitando duplicidade de pagador. Um
cadastro inativo com o mesmo CPF não impede a inclusão.

\fluxo{E2}{Inativação de aluno com cobrança em aberto}
Se o aluno possuir cobrança com situação ``Aberta'' ou ``Vencida'', o sistema informa a existência
do débito e exige confirmação explícita (Figura~\ref{fig:aluno-inativar}). As cobranças em aberto
não são canceladas automaticamente: permanecem válidas e continuam a compor a inadimplência
([Indicadores.RF002]).

\textit{Observação sobre o protótipo:} os formulários das Figuras~\ref{fig:cadastros-novo-aluno}
e~\ref{fig:aluno-dados} ainda apresentam os campos ``Data de nascimento'' e ``Telefone'', anteriores
à definição do modelo de dados. Prevalecem as entradas descritas neste caso de uso.

\tela{cadastros-alunos}{Cadastros, aba Alunos: busca por nome ou CPF, filtro por situação, valor
em aberto e sinalização de cadastro incompleto.}{cadastros-alunos}
\tela{cadastros-dialogo-novo-aluno}{Inclusão de aluno, com o aviso de CPF já cadastrado
(fluxo E1).}{cadastros-novo-aluno}
\tela{aluno-dados}{Cadastro do aluno, aba Dados: alteração dos dados cadastrais.}{aluno-dados}
\tela{aluno-cobrancas}{Cadastro do aluno, aba Cobranças: cobranças em seu nome e total em
aberto.}{aluno-cobrancas}
\tela{aluno-movimentacoes}{Cadastro do aluno, aba Movimentações: entradas geradas pelas cobranças
pagas.}{aluno-movimentacoes}
\tela{aluno-dialogo-inativar}{Inativação de aluno com débito: resumo do efeito e confirmação
(fluxo E2).}{aluno-inativar}

\FloatBarrier

% ---------------------------------------------------------------------------
\req{RF002}{Manter cadastro de funcionários}

Permite incluir, alterar, consultar e inativar o cadastro dos funcionários do CASD, aos quais são
vinculadas as movimentações de saída referentes a pagamento de pessoal.

\prioE

\entradas o usuário está autenticado. Entradas: nome completo, CPF, cargo, salário de referência,
data de admissão, data de desligamento quando houver e situação.

\saidas o cadastro é criado, alterado ou inativado, e passa a estar disponível para vinculação em
[Movimentações.RF001]. O salário de referência é informativo: o pagamento efetivo é sempre lançado
como movimentação de saída vinculada ao funcionário.

\fluxoprincipal
\begin{enumerate}
  \item O usuário seleciona a opção ``Cadastros'', aba ``Funcionários'', e pesquisa pelo nome ou
        CPF (Figura~\ref{fig:cadastros-funcionarios}).
  \item O sistema apresenta o cadastro correspondente ou oferece a inclusão de um novo
        funcionário.
  \item O usuário preenche ou altera os dados.
  \item O sistema valida os campos obrigatórios e grava o cadastro.
\end{enumerate}

\fluxossecundarios

\fluxo{A1}{Desligamento}
O usuário informa a data de desligamento e o sistema inativa o funcionário automaticamente a partir
dessa data, mantendo o histórico de pagamentos associado.

\fluxo{E1}{CPF já cadastrado}
Comportamento idêntico ao descrito em [Cadastros.RF001], fluxo E1.

\tela{cadastros-funcionarios}{Cadastros, aba Funcionários: cargo, admissão, desligamento e
situação.}{cadastros-funcionarios}

\FloatBarrier

% ---------------------------------------------------------------------------
\reqretirado{RF003}{Manter cadastro de fornecedores e parceiros}{Retirado do escopo do MVP na
versão 1.1. O contrato registra a contraparte apenas pelo nome ([Cadastros.RF004]), e as
movimentações deixam de ser vinculadas a fornecedor ([Movimentações.RF001]). O identificador
permanece reservado.}

% ---------------------------------------------------------------------------
\req{RF004}{Manter contratos}

Permite registrar e acompanhar os contratos firmados pelo CASD, que são a origem das despesas
recorrentes do Centro Acadêmico.

\prioI

\entradas o usuário está autenticado. Entradas: nome da contraparte, objeto do contrato, valor,
periodicidade do pagamento (mensal, por demanda ou única), data de início e data de término.

\saidas o contrato é registrado e fica disponível para vinculação nas movimentações de saída,
permitindo apurar quanto foi efetivamente pago em cada contrato.

\fluxoprincipal
\begin{enumerate}
  \item O usuário seleciona a opção ``Cadastros'', aba ``Contratos''
        (Figura~\ref{fig:cadastros-contratos}).
  \item O sistema apresenta a relação de contratos, com a vigência, o total pago no ano corrente e a
        situação de cada um, destacando os que vencem nos próximos sessenta dias.
  \item O usuário inclui um novo contrato ou seleciona um existente para alteração.
  \item O usuário informa a contraparte, o objeto, o valor, a periodicidade e a vigência.
  \item O sistema valida os dados e grava o contrato.
\end{enumerate}

\fluxossecundarios

\fluxo{A1}{Encerramento antecipado}
O usuário informa a data efetiva de encerramento e a justificativa. O contrato passa à situação
``Encerrado'' e deixa de aparecer entre os vigentes, sem afetar as movimentações já registradas.

\fluxo{E1}{Vigência inconsistente}
Se a data de término for anterior à de início, o sistema recusa a gravação e indica o campo
inconsistente.

\tela{cadastros-contratos}{Cadastros, aba Contratos: aviso de vencimento próximo, contraparte em
texto e total pago no ano.}{cadastros-contratos}

\FloatBarrier

% ---------------------------------------------------------------------------
\req{RF005}{Manter projetos}

Permite registrar os projetos do CASD, que funcionam como unidade de agrupamento de despesas e
receitas e permitem responder quanto custou cada iniciativa do Centro Acadêmico.

\prioI

\entradas o usuário está autenticado. Entradas: nome do projeto, responsável, período de execução,
orçamento previsto quando houver e situação.

\saidas o projeto fica disponível para vinculação em [Movimentações.RF001] e passa a compor os
recortes de consulta das movimentações.

\fluxoprincipal
\begin{enumerate}
  \item O usuário seleciona a opção ``Cadastros'', aba ``Projetos''
        (Figura~\ref{fig:cadastros-projetos}).
  \item O sistema apresenta a relação de projetos, separando os em andamento dos encerrados, com o
        total movimentado e o orçamento previsto de cada um.
  \item O usuário inclui um novo projeto ou seleciona um existente.
  \item O usuário informa nome, responsável, período e, opcionalmente, orçamento previsto.
  \item O sistema valida os dados e grava o projeto.
\end{enumerate}

\fluxossecundarios

\fluxo{A1}{Encerramento de projeto}
O usuário encerra o projeto e o sistema passa a exibir, junto ao registro, o total efetivamente
movimentado e a comparação com o orçamento previsto, quando este tiver sido informado.

\fluxo{E1}{Encerramento com movimentação posterior}
Movimentações podem ser vinculadas a projeto encerrado, já que despesas costumam chegar depois do
término. O sistema aceita o lançamento e sinaliza a situação ao usuário, em vez de recusá-lo.

\tela{cadastros-projetos}{Cadastros, aba Projetos: total movimentado comparado ao orçamento
previsto.}{cadastros-projetos}

\FloatBarrier

% ===========================================================================
\section{Cobranças}

Esta subseção reúne todo o ciclo de vida da cobrança dirigida ao aluno, da emissão à baixa. Os
casos de uso operam sobre a mesma entidade e compartilham as regras de integração com o Asaas, que
é a plataforma responsável por emitir o boleto, entregá-lo ao pagador e receber o pagamento.

Três regras gerais valem para toda a subseção. A primeira é que o IntegraCASD não recebe dinheiro:
ele solicita a cobrança ao Asaas e registra o desfecho informado pela plataforma. A segunda é que
nenhuma cobrança existe apenas no IntegraCASD: se o Asaas falhar, nada é gravado
([Confiabilidade.NF005]). A terceira é que multa e juros de atraso são aplicados pelo Asaas, com os
percentuais enviados na criação de cada cobrança ([Administração.RF003]); o sistema não calcula
encargos, e o valor que exibe é sempre o valor original da cobrança.

A lista de cobranças (Figura~\ref{fig:cobrancas}) apenas mostra e leva ao detalhe. O detalhe da
cobrança (Figura~\ref{fig:cobranca}) é o centro das ações: dele saem a obtenção do boleto, a
sincronização, o registro manual de pagamento e o cancelamento.

Os atores variam entre os casos de uso e estão indicados em cada bloco.

% ---------------------------------------------------------------------------
\req{RF001}{Gerar cobranças recorrentes}

Cria automaticamente, a cada período, as cobranças de mensalidade dos alunos ativos, sem intervenção
humana.

\ator{Tempo.}
\prioE

\entradas existem parâmetros de cobrança recorrente configurados ([Administração.RF003]); existe
ao menos um aluno ativo; a integração com o Asaas está configurada e operacional
([Administração.RF004]). Entradas: data corrente e parâmetros vigentes.

\saidas para cada aluno ativo é criada uma cobrança com situação ``Aberta'', associada a uma
cobrança correspondente no Asaas, com valor, vencimento e percentuais de multa e juros definidos
pelos parâmetros. O Asaas envia o boleto por e-mail ao aluno. O resultado da execução fica
registrado, com o total de cobranças geradas e as que falharam, e pode ser consultado na tela de
parâmetros (Figura~\ref{fig:adm-parametros}).

\fluxoprincipal
\begin{enumerate}
  \item Ao chegar o dia de geração definido nos parâmetros, o sistema inicia a rotina.
  \item O sistema obtém a relação de alunos ativos.
  \item Para cada aluno, o sistema verifica se já existe cobrança de mensalidade para o período de
        referência.
  \item Não havendo, o sistema solicita ao Asaas a criação da cobrança, informando pagador, valor,
        data de vencimento e percentuais de multa e juros.
  \item O sistema registra a cobrança retornada pelo Asaas, com a sua identificação na plataforma e
        a situação ``Aberta''.
  \item Ao final, o sistema grava o resultado da execução, com o horário de início e de término, o
        número de cobranças criadas e a relação das que falharam.
\end{enumerate}

\fluxossecundarios

\fluxo{A1}{Cobrança já existente para o período}
Se já houver cobrança de mensalidade para aquele aluno e período de referência, o sistema não cria
outra e segue para o próximo aluno. Essa verificação torna a rotina segura contra execução
repetida, o que importa porque uma falha parcial exige nova execução.

\fluxo{A2}{Aluno sem e-mail ou sem CPF válido}
O sistema não cria a cobrança, registra a pendência na relação de falhas com a identificação do
aluno e do campo faltante, e prossegue com os demais. A área administrativa corrige o cadastro e
executa novamente a geração para os pendentes.

\fluxo{E1}{Asaas indisponível}
Se a plataforma não responder, o sistema interrompe a rotina, registra quais cobranças já haviam
sido criadas e agenda nova tentativa. Nenhuma cobrança é criada apenas localmente.

\fluxo{E2}{Falha após criar a cobrança no Asaas}
Se a cobrança for criada na plataforma mas o registro local falhar, a sincronização
([Cobranças.RF007]) identifica a divergência e registra a cobrança faltante, em vez de criar uma
duplicata.

% ---------------------------------------------------------------------------
\req{RF002}{Emitir cobrança avulsa}

Permite emitir, em uma única operação, uma cobrança fora do ciclo recorrente, dirigida a um aluno
específico, como a inscrição em um evento ou a venda de um produto do CASD.

\ator{Área administrativa.}
\prioE

\entradas o aluno está cadastrado e ativo, com e-mail e CPF válidos; a integração com o Asaas está
operacional. Entradas: aluno, valor, data de vencimento e descrição da cobrança.

\saidas a cobrança é criada no Asaas e registrada no sistema com situação ``Aberta''; o boleto é
enviado por e-mail ao aluno pela plataforma; a confirmação da emissão fica disponível em tela.

\fluxoprincipal
\begin{enumerate}
  \item O usuário aciona ``Nova cobrança'', disponível na barra superior de todas as telas e no
        cadastro do aluno.
  \item O sistema apresenta o formulário em uma janela sobre a tela atual
        (Figura~\ref{fig:nova-cobranca}), com data de vencimento sugerida a partir dos parâmetros
        vigentes e, quando aberto a partir do cadastro do aluno, com o aluno já preenchido.
  \item O usuário seleciona o aluno por nome ou CPF e informa valor, vencimento e descrição.
  \item O usuário confirma a emissão.
  \item O sistema solicita a criação da cobrança ao Asaas, informando também os percentuais de
        multa e juros.
  \item O sistema registra a cobrança retornada e exibe a confirmação, com a opção de obter o
        boleto ([Cobranças.RF003]).
\end{enumerate}

\fluxossecundarios

\fluxo{A1}{Aluno com cadastro incompleto}
Se faltar e-mail ou CPF, o sistema informa o campo pendente e oferece atalho para o cadastro do
aluno ([Cadastros.RF001]), retornando ao formulário depois da correção.

\fluxo{E1}{Asaas indisponível}
O sistema informa a indisponibilidade, não registra a cobrança localmente e orienta nova tentativa.
A operação não é enfileirada para execução posterior, para que a área administrativa saiba com
certeza se a cobrança existe.

\fluxo{E2}{Vencimento no passado}
Se a data de vencimento informada for anterior à data corrente, o sistema recusa a emissão e
solicita correção.

\tela{topo-dialogo-nova-cobranca}{Nova cobrança avulsa, aberta pela barra superior.}{nova-cobranca}

\FloatBarrier

% ---------------------------------------------------------------------------
\req{RF003}{Obter o boleto de uma cobrança}

Permite baixar o boleto de uma cobrança, para que a área administrativa o entregue ao aluno.

\ator{Área administrativa.}
\prioI

\entradas existe uma cobrança com situação ``Aberta'' ou ``Vencida''. Entrada: identificação da
cobrança.

\saidas o arquivo do boleto é disponibilizado para download, e a operação é registrada no
histórico da cobrança, com autor, data e hora.

\fluxoprincipal
\begin{enumerate}
  \item O usuário localiza a cobrança ([Cobranças.RF008]) e abre o seu detalhe
        (Figura~\ref{fig:cobranca}).
  \item O sistema apresenta os dados da cobrança, a data da última sincronização com o Asaas, o
        histórico de eventos e as ações disponíveis.
  \item O usuário aciona ``Baixar boleto''.
  \item O sistema obtém o boleto no Asaas e o disponibiliza para download.
  \item O sistema registra a operação.
\end{enumerate}

\fluxossecundarios

\fluxo{A1}{Entrega por outro meio}
Ao optar pelo download, o usuário obtém o arquivo e passa a ser responsável pela entrega, que pode
ocorrer por qualquer canal. O sistema não acompanha essa entrega.

\fluxo{E1}{Cobrança já paga ou cancelada}
O sistema não disponibiliza o boleto e informa a situação atual da cobrança, evitando que um
documento sem validade circule.

O reenvio do boleto por e-mail, previsto até a versão 1.1, foi retirado do MVP na versão 1.2 (ver
o Histórico de revisões).

\tela[0.5]{cobranca}{Detalhe da cobrança: ações, dados, situação da sincronização e histórico de
eventos do Asaas.}{cobranca}

\FloatBarrier

% ---------------------------------------------------------------------------
\req{RF004}{Cancelar cobrança}

Permite cancelar uma cobrança ainda não paga, tanto no IntegraCASD quanto no Asaas. É também o
mecanismo pelo qual se opera uma renegociação, que consiste em cancelar a cobrança vigente e emitir
outra em condições novas.

\ator{Área administrativa.}
\prioI

\entradas existe uma cobrança com situação ``Aberta'' ou ``Vencida''. Entradas: identificação da
cobrança e motivo do cancelamento.

\saidas a cobrança passa à situação ``Cancelada'' nos dois sistemas, deixa de compor a
inadimplência e o seu boleto perde a validade. O motivo, o autor, a data e a hora ficam
registrados.

\fluxoprincipal
\begin{enumerate}
  \item No detalhe da cobrança, o usuário aciona ``Cancelar / renegociar''.
  \item O sistema apresenta o resumo do efeito da operação e solicita o motivo
        (Figura~\ref{fig:cobranca-cancelar}).
  \item O usuário informa o motivo e confirma.
  \item O sistema solicita o cancelamento ao Asaas.
  \item Confirmado o cancelamento na plataforma, o sistema atualiza a situação da cobrança e
        registra a operação.
\end{enumerate}

\fluxossecundarios

\fluxo{A1}{Renegociação}
Se o usuário marcar a opção de emitir nova cobrança em seguida, o sistema, concluído o
cancelamento, abre a emissão de cobrança avulsa ao mesmo aluno ([Cobranças.RF002]) com os dados da
cobrança cancelada pré-preenchidos. A vinculação entre as duas fica registrada, de modo que o
histórico do aluno mostre que se trata de renegociação, e não de duas dívidas distintas.

\fluxo{E1}{Cobrança já paga}
Se a cobrança já estiver paga, o sistema recusa o cancelamento e informa que a devolução de valores
é feita pelo Asaas, cujo estorno chega ao sistema pela conciliação ([Cobranças.RF005]).

\fluxo{E2}{Falha no cancelamento junto ao Asaas}
Se a plataforma recusar ou não responder, o sistema não altera a situação local e informa a falha.
Cancelar apenas no IntegraCASD deixaria um boleto válido em circulação, com risco de o aluno pagar
uma cobrança que o CASD considera inexistente.

\tela{cobranca-dialogo-cancelar}{Cancelamento de cobrança: resumo do efeito, motivo e opção de
renegociação.}{cobranca-cancelar}

\FloatBarrier

% ---------------------------------------------------------------------------
\req{RF005}{Conciliar pagamento de cobrança}

Atualiza automaticamente a situação de uma cobrança a partir das notificações enviadas pelo Asaas,
sem intervenção humana. É o mecanismo que dispensa a conferência manual contra o extrato bancário.

\ator{Asaas.}
\prioE

\entradas a integração está configurada e o endereço de notificação registrado na plataforma.
Entrada: notificação de evento de cobrança enviada pelo Asaas.

\saidas a situação da cobrança é atualizada; havendo pagamento, é criada a movimentação de entrada
correspondente, vinculada à cobrança; a notificação processada fica registrada, com data, hora e
conteúdo, e aparece no histórico de eventos da cobrança (Figura~\ref{fig:cobranca}).

\fluxoprincipal
\begin{enumerate}
  \item O Asaas envia ao sistema a notificação de um evento de cobrança.
  \item O sistema verifica a autenticidade da notificação pelo token \texttt{asaas-access-token}.
  \item O sistema grava a notificação, enfileira o seu processamento e confirma o recebimento ao
        Asaas, sem processá-la nesse momento ([Desempenho.NF001]).
  \item Em seguida, de forma desacoplada, o sistema localiza a cobrança correspondente pela
        identificação da plataforma.
  \item O sistema atualiza a situação da cobrança conforme o evento recebido: pagamento confirmado
        e pagamento recebido levam à situação ``Paga''; vencimento sem pagamento leva à situação
        ``Vencida''; estorno leva à situação ``Estornada''.
  \item Tratando-se de pagamento, o sistema cria a movimentação de entrada correspondente, com a
        categoria associada ao tipo da cobrança e forma de pagamento ``Asaas'', na mesma transação
        em que atualiza a cobrança ([Confiabilidade.NF004]).
  \item O sistema marca a notificação como processada.
\end{enumerate}

\fluxossecundarios

\fluxo{A1}{Distinção entre pagamento confirmado e valor disponível}
O Asaas sinaliza separadamente o pagamento processado e a disponibilidade do valor em conta. O
sistema considera a cobrança paga já na confirmação e registra a data de disponibilidade quando ela
chega, de modo que o painel possa distinguir o resultado do período do valor ainda aguardando
liberação.

\fluxo{A2}{Notificação repetida}
Se a notificação já tiver sido processada, identificada pelo seu identificador no Asaas, o sistema
a descarta sem alterar nada e confirma o recebimento. Sem essa verificação, um reenvio da
plataforma criaria uma segunda movimentação de entrada para o mesmo pagamento
([Confiabilidade.NF003]).

\fluxo{E1}{Cobrança não encontrada}
Se a notificação referir-se a cobrança inexistente no sistema, ela é registrada como pendência de
conciliação e apresentada à área administrativa, que pode resolvê-la pela sincronização
([Cobranças.RF007]). A notificação não é descartada.

\fluxo{E2}{Notificação não autêntica}
Notificação cuja origem não possa ser verificada é rejeitada e registrada como tentativa de acesso
indevido, sem produzir qualquer alteração. Como o endereço de notificação é público, essa
verificação é o que impede a alteração de dados financeiros por terceiros.

% ---------------------------------------------------------------------------
\req{RF006}{Registrar pagamento manualmente}

Permite que a área administrativa registre o pagamento de uma cobrança que não foi refletido pela
conciliação automática: seja porque o valor comprovadamente entrou e a notificação não chegou, seja
porque o pagamento foi combinado e recebido diretamente pela área administrativa, por exemplo por
Pix na conta do CASD. É um mecanismo de exceção, e não um fluxo de rotina.

\ator{Área administrativa.}
\prioI

\entradas existe cobrança com situação ``Aberta'' ou ``Vencida''. Entradas: identificação da
cobrança, data do pagamento, valor pago e justificativa escrita.

\saidas a cobrança passa à situação ``Paga com registro manual'', distinta de ``Paga'', de modo que
a origem do dado permaneça visível; a movimentação de entrada correspondente é criada; a
justificativa, o autor, a data e a hora ficam registrados.

\fluxoprincipal
\begin{enumerate}
  \item No detalhe da cobrança, o usuário aciona ``Registrar pagamento manualmente''.
  \item O sistema apresenta a data da última sincronização da cobrança com o Asaas e recomenda
        sincronizar antes, porque o pagamento pode já estar na plataforma
        (Figura~\ref{fig:cobranca-pagamento-manual}).
  \item O usuário informa data do pagamento, valor e justificativa.
  \item O sistema apresenta o resumo do efeito, alertando que o registro fará o sistema divergir da
        plataforma, e exige confirmação.
  \item O usuário confirma.
  \item O sistema atualiza a situação da cobrança, cria a movimentação de entrada e registra a
        operação com a justificativa, na mesma transação ([Confiabilidade.NF004]).
\end{enumerate}

\fluxossecundarios

\fluxo{A1}{Sincronização resolve a divergência}
Se o usuário sincronizar a cobrança antes de prosseguir e o Asaas indicar o pagamento, a cobrança
é baixada pela sincronização ([Cobranças.RF007]) e o caso de uso é encerrado sem registro manual.

\fluxo{E1}{Justificativa não informada}
O sistema não conclui a operação sem justificativa escrita. É o único registro de por que o sistema
passou a divergir da plataforma.

\fluxo{E2}{Conciliação posterior do mesmo pagamento}
Se, depois do registro manual, chegar do Asaas a notificação de pagamento daquela mesma cobrança, o
sistema não cria segunda movimentação: vincula a notificação ao registro manual e sinaliza a
ocorrência para conferência, já que pode indicar pagamento em duplicidade pelo aluno.

\tela{cobranca-dialogo-pagamento-manual}{Registro manual de pagamento: alerta de sincronização
antiga, justificativa obrigatória e resumo do efeito.}{cobranca-pagamento-manual}

\FloatBarrier

% ---------------------------------------------------------------------------
\req{RF007}{Sincronizar cobranças com o Asaas}

Consulta a plataforma e atualiza a situação das cobranças, recuperando eventos que não tenham
chegado por notificação.

\ator{Área administrativa e Tempo.}
\prioI

\entradas a integração está configurada e operacional. Entrada: período de vencimento a
sincronizar, ou uma cobrança específica.

\saidas as cobranças cuja situação divergia são atualizadas; as movimentações de entrada
correspondentes aos pagamentos recuperados são criadas; um resumo do que foi alterado é
apresentado ao usuário e registrado, com a data e a hora da sincronização.

\fluxoprincipal
\begin{enumerate}
  \item Na lista de cobranças, o usuário aciona ``Sincronizar com o Asaas'' e informa o período
        (Figura~\ref{fig:cobrancas-sincronizar}).
  \item O sistema consulta na plataforma a situação das cobranças daquele período.
  \item O sistema compara a situação de cada cobrança com a registrada localmente.
  \item Havendo divergência, o sistema adota a situação informada pela plataforma e cria as
        movimentações necessárias.
  \item O sistema apresenta o resumo do que foi alterado (cobranças conferidas, situações
        atualizadas, cobranças encontradas só no Asaas e cobranças com registro manual a conferir)
        e registra a execução.
\end{enumerate}

\fluxossecundarios

\fluxo{A1}{Sincronização de uma única cobrança}
A partir do detalhe da cobrança, o usuário pode sincronizar apenas aquela cobrança
(Figura~\ref{fig:cobranca-sincronizar}). É o passo recomendado antes de cobrar o aluno ou de
registrar pagamento manualmente.

\fluxo{A2}{Execução periódica}
O sistema executa a sincronização automaticamente uma vez por dia, sem intervenção humana, e
registra o resultado como as demais rotinas automáticas ([Confiabilidade.NF007]).

\fluxo{A3}{Cobrança existente apenas no Asaas}
Cobrança que exista na plataforma e não no sistema é registrada localmente, o que trata a falha
descrita em [Cobranças.RF001], fluxo E2.

\fluxo{A4}{Cobrança registrada manualmente}
Cobrança com situação ``Paga com registro manual'' não é sobrescrita automaticamente. A divergência
é apresentada ao usuário para conferência, porque sobrescrevê-la apagaria a justificativa que a
originou.

\fluxo{E1}{Asaas indisponível}
O sistema informa a indisponibilidade e não altera nada, mantendo as situações locais.

\tela{cobrancas-dialogo-sincronizar}{Sincronização geral a partir da lista de cobranças, com o
resultado da última execução.}{cobrancas-sincronizar}
\tela{cobranca-dialogo-sincronizar}{Sincronização de uma única cobrança, a partir do seu detalhe
(fluxo A1).}{cobranca-sincronizar}

\FloatBarrier

% ---------------------------------------------------------------------------
\req{RF008}{Consultar cobranças}

Permite localizar cobranças e acompanhar a sua situação.

\ator{Área administrativa.}
\prioE

\entradas o usuário está autenticado. Entradas: filtros por aluno, situação, período de vencimento
e faixa de valor.

\saidas a relação de cobranças é apresentada, com aluno, descrição, valor, vencimento, situação e
dias de atraso quando houver. O sistema não sofre alteração de estado.

\fluxoprincipal
\begin{enumerate}
  \item O usuário seleciona a opção ``Cobranças'' (Figura~\ref{fig:cobrancas}).
  \item O sistema apresenta o estado da conciliação (hora da última notificação e da última
        sincronização), os totais por situação e, por padrão, as cobranças abertas e vencidas,
        ordenadas por vencimento.
  \item O usuário ajusta os filtros ou seleciona outra situação pelos cartões de resumo.
  \item O sistema apresenta o resultado.
  \item O usuário pode selecionar uma cobrança para ver o detalhe e acessar as operações a ela
        aplicáveis.
\end{enumerate}

\fluxossecundarios

\fluxo{A1}{Consulta a partir do histórico do aluno}
A consulta também pode ser iniciada a partir do cadastro do aluno, hipótese em que o filtro já vem
aplicado ([Cadastros.RF001], fluxo A3).

\fluxo{A2}{Sincronização antiga}
Cobrança vencida cuja última sincronização com o Asaas seja antiga aparece sinalizada na lista,
porque o aluno pode ter pago e a notificação não ter chegado.

\tela[0.5]{cobrancas}{Lista de cobranças: estado da conciliação, totais por situação, filtros e
cobranças abertas e vencidas.}{cobrancas}

\FloatBarrier

% ===========================================================================
\section{Movimentações}

Esta subseção reúne o registro e a consulta de todas as entradas e saídas de recursos do CASD. As
movimentações têm duas origens: são criadas automaticamente pela conciliação de cobranças
([Cobranças.RF005]) e pelo registro manual de pagamento ([Cobranças.RF006]), ou lançadas
manualmente pela área administrativa, no caso de pagamentos de contratos, despesas de projeto,
pessoal e recebimentos avulsos, que ocorrem por Pix, transferência, dinheiro ou cartão.

Vale para toda a subseção a regra de que movimentação registrada não é editada nem excluída.
Correções são feitas por estorno, na forma de [Movimentações.RF002].

\textbf{Ator de todos os casos de uso desta subseção:} área administrativa.

% ---------------------------------------------------------------------------
\req{RF001}{Registrar movimentação}

Permite lançar uma entrada ou uma saída de recursos, classificando-a e vinculando-a aos cadastros a
que se refere.

\prioE

\entradas existe ao menos uma categoria de movimentação ativa ([Administração.RF002]). Entradas:
tipo, valor, data, categoria, forma de pagamento (Pix, transferência, dinheiro ou cartão),
descrição e, opcionalmente, vínculo com um contrato, um projeto ou um funcionário.

\saidas a movimentação é registrada e passa a compor o resultado do período e os indicadores. Fica
registrado o autor, com data e hora.

\fluxoprincipal
\begin{enumerate}
  \item O usuário aciona ``Nova movimentação'', disponível na barra superior de todas as telas.
  \item O sistema apresenta o formulário em uma janela sobre a tela atual
        (Figura~\ref{fig:nova-movimentacao}), com a data corrente sugerida.
  \item O usuário informa o tipo, o valor, a data, a categoria e a forma de pagamento.
  \item O usuário informa a descrição e, quando aplicável, vincula a movimentação a contrato,
        projeto ou funcionário.
  \item O sistema valida os campos obrigatórios e grava a movimentação.
\end{enumerate}

\fluxossecundarios

\fluxo{A1}{Movimentação vinculada a contrato}
Ao vincular a movimentação a um contrato, o sistema sugere o valor previsto e apresenta o total já
pago naquele contrato no ano, permitindo que o usuário perceba divergências no momento do
lançamento.

\fluxo{A2}{Movimentação vinculada a projeto encerrado}
O sistema aceita o lançamento e sinaliza a situação, conforme [Cadastros.RF005], fluxo E1.

\fluxo{E1}{Data futura}
O sistema recusa o lançamento com data posterior à corrente, porque movimentação registra fato
ocorrido, e não previsão.

\fluxo{E2}{Valor igual a zero}
O sistema recusa a gravação e solicita correção.

\tela{topo-dialogo-nova-movimentacao}{Nova movimentação, aberta pela barra superior, com a
sugestão de valor do contrato vinculado (fluxo A1).}{nova-movimentacao}

\FloatBarrier

% ---------------------------------------------------------------------------
\req{RF002}{Estornar movimentação}

Permite corrigir uma movimentação errada sem apagar o registro original, por meio de um lançamento
de sinal contrário que anula o seu efeito.

\prioI

\entradas existe movimentação lançada manualmente e ainda não estornada. Entradas: identificação
da movimentação e motivo do estorno.

\saidas a movimentação original permanece no sistema, marcada como estornada; é criada uma
movimentação espelho, de sinal contrário, vinculada à original; o efeito das duas sobre o resultado
é nulo; o motivo, o autor, a data e a hora ficam registrados.

\fluxoprincipal
\begin{enumerate}
  \item Na lista de movimentações, o usuário aciona ``Estornar'' na movimentação desejada.
  \item O sistema apresenta os dados da movimentação, o resumo do efeito e solicita o motivo
        (Figura~\ref{fig:movimentacoes-estornar}).
  \item O usuário informa o motivo e confirma.
  \item O sistema marca a movimentação original como estornada e cria a movimentação espelho.
  \item O sistema registra a operação e, se o usuário tiver marcado a opção correspondente, abre o
        lançamento da movimentação correta.
\end{enumerate}

\fluxossecundarios

\fluxo{E1}{Movimentação originada de cobrança}
O sistema não oferece o estorno de movimentação criada pela conciliação ([Cobranças.RF005]) ou pelo
registro manual de pagamento ([Cobranças.RF006]) e orienta o usuário a tratar a cobrança
correspondente. Essas movimentações são espelho do que ocorreu no Asaas, e estorná-las diretamente
faria o sistema divergir da plataforma. Devolução ao aluno gera evento de estorno na plataforma,
que chega pela conciliação.

\fluxo{E2}{Movimentação já estornada}
O sistema informa que a movimentação já foi estornada e apresenta o lançamento espelho
correspondente, sem criar outro.

\tela{movimentacoes-dialogo-estornar}{Estorno de movimentação: resumo do efeito, motivo e opção de
lançar a movimentação correta.}{movimentacoes-estornar}

\FloatBarrier

% ---------------------------------------------------------------------------
\req{RF003}{Consultar movimentações}

Permite localizar movimentações e apurar totais por recorte.

\prioE

\entradas o usuário está autenticado. Entradas: filtros por período, tipo, categoria, forma de
pagamento, projeto ou contrato.

\saidas a relação de movimentações é apresentada, com data, descrição, vínculo, categoria, forma,
origem (manual, conciliação automática, registro manual de cobrança ou estorno) e valor, além dos
totais de entrada, de saída e do resultado do recorte. Movimentações estornadas aparecem
sinalizadas, junto ao seu espelho. O sistema não sofre alteração de estado.

\fluxoprincipal
\begin{enumerate}
  \item O usuário seleciona a opção ``Movimentações'' (Figura~\ref{fig:movimentacoes}).
  \item O sistema apresenta, por padrão, as movimentações do mês corrente, com os totais.
  \item O usuário ajusta os filtros.
  \item O sistema apresenta o resultado com os totais do recorte.
\end{enumerate}

\fluxossecundarios

\fluxo{A1}{Ocultar estornos}
O usuário pode optar por ocultar as movimentações estornadas e os seus espelhos. Os totais não
mudam, já que o par se anula, mas a leitura da relação fica mais simples.

\fluxo{A2}{Consulta a partir de um indicador}
A consulta também pode ser aberta a partir do painel de indicadores, hipótese em que o filtro do
indicador selecionado já vem aplicado ([Indicadores.RF001], [Indicadores.RF003]).

\tela[0.5]{movimentacoes}{Lista de movimentações: totais do recorte, filtros, origem de cada
lançamento e movimentação estornada com o seu espelho.}{movimentacoes}

\FloatBarrier

% ---------------------------------------------------------------------------
\reqretirado{RF004}{Anexar comprovante a uma movimentação}{Retirado do escopo do MVP na versão
1.1. O identificador permanece reservado.}

% ===========================================================================
\section{Indicadores}

Esta subseção reúne as consultas que apoiam a decisão. Os dois perfis de usuário acessam os mesmos
indicadores, com finalidades distintas: a área administrativa os utiliza para agir, sobretudo sobre
a inadimplência, e a diretoria os utiliza para decidir sobre gastos, contratos e projetos. Os três
indicadores ficam reunidos em uma única tela, que é a página inicial do sistema
(Figura~\ref{fig:indicadores}).

\textbf{Ator de todos os casos de uso desta subseção:} área administrativa e diretoria.

% ---------------------------------------------------------------------------
\req{RF001}{Consultar painel de indicadores}

Apresenta a posição financeira do CASD em um período, reunindo os números que descrevem entradas,
saídas e resultado.

\prioE

\entradas existem movimentações registradas. Entradas: período de referência, com o mês corrente
como padrão.

\saidas o painel apresenta o total de entradas, o total de saídas e o resultado do período,
comparados ao período anterior, o valor já pago pelos alunos e ainda aguardando liberação pelo
Asaas, e a evolução de entradas e saídas nos seis meses anteriores, além dos indicadores de
inadimplência ([Indicadores.RF002]) e da distribuição por categoria ([Indicadores.RF003]). O
sistema não sofre alteração de estado.

\fluxoprincipal
\begin{enumerate}
  \item O usuário acessa o sistema e o painel é apresentado como página inicial.
  \item O sistema apura os indicadores do mês corrente e os exibe.
  \item O usuário pode alterar o período de referência.
  \item O sistema recalcula e apresenta os indicadores do novo período.
  \item O usuário pode selecionar um indicador para ver as movimentações que o compõem.
\end{enumerate}

\fluxossecundarios

\fluxo{A1}{Período sem movimentação}
O sistema apresenta o painel com valores zerados e a informação de que não há movimentação no
período, em vez de retornar erro.

\fluxo{A2}{Distinção entre resultado e disponibilidade}
Quando houver pagamentos confirmados cujo valor ainda não esteja disponível em conta, o painel
indica os dois números separadamente, de modo que o resultado apresentado não seja confundido com
dinheiro disponível.

\fluxo{A3}{Dados do gráfico em tabela}
Os valores do gráfico de evolução também podem ser consultados em forma de tabela
([Usabilidade.NF004]).

% ---------------------------------------------------------------------------
\req{RF002}{Consultar inadimplência}

Apresenta a situação dos alunos com cobranças vencidas e não pagas.

\prioE

\entradas existem cobranças registradas. Entradas: período de referência.

\saidas são apresentados o número de alunos inadimplentes, o valor total em aberto, a distribuição
dos inadimplentes por tempo de atraso (até 30 dias, de 31 a 60 dias e mais de 60 dias) e a relação
dos maiores devedores, com o número de cobranças e o atraso da mais antiga. Os valores exibidos são
os originais das cobranças, sem multa e juros, que são apurados pelo Asaas apenas no momento do
pagamento. O sistema não sofre alteração de estado.

\fluxoprincipal
\begin{enumerate}
  \item O usuário acessa o indicador de inadimplência, no painel.
  \item O sistema apura as cobranças com situação ``Vencida'' e as agrupa por aluno.
  \item O sistema apresenta o total de inadimplentes, o valor em aberto e a relação de devedores,
        ordenada por valor.
  \item O usuário pode selecionar um aluno para ver as suas cobranças vencidas, ou abrir a lista
        completa de cobranças vencidas ([Cobranças.RF008]).
\end{enumerate}

\fluxossecundarios

\fluxo{A1}{Ação de cobrança a partir da relação}
A partir da relação, o usuário com perfil administrativo pode abrir as cobranças vencidas de um
aluno e obter o boleto ([Cobranças.RF003]), para entregá-lo ao aluno, o que atende ao uso que a
área administrativa faz deste indicador.

\fluxo{A2}{Divergência suspeita}
Se houver cobranças vencidas cuja última sincronização com o Asaas tenha ocorrido há mais de sete
dias, o sistema sinaliza a situação no topo do painel e oferece a sincronização
([Cobranças.RF007]), reduzindo a chance de um aluno figurar como inadimplente por falha de
integração.

% ---------------------------------------------------------------------------
\req{RF003}{Consultar gastos por categoria}

Apresenta a distribuição das saídas por categoria de movimentação, permitindo identificar em que o
CASD gasta os seus recursos.

\prioI

\entradas existem movimentações de saída no período. Entradas: período de referência.

\saidas é apresentada a distribuição das saídas por categoria, com valor e participação
percentual, ordenada por valor. O sistema não sofre alteração de estado.

\fluxoprincipal
\begin{enumerate}
  \item O usuário acessa o indicador de gastos por categoria, no painel.
  \item O sistema apura as saídas do período agrupadas por categoria.
  \item O sistema apresenta a distribuição, com valores e percentuais.
  \item O usuário pode selecionar uma categoria para ver as movimentações que a compõem
        ([Movimentações.RF003]).
\end{enumerate}

\fluxossecundarios

\fluxo{A1}{Comparação com período anterior}
O usuário pode solicitar a comparação com o período imediatamente anterior, hipótese em que o
sistema apresenta a variação de cada categoria.

\fluxo{A2}{Movimentação sem categoria}
Movimentações que porventura estejam sem categoria, como as vindas da carga inicial, são agrupadas
em ``Não classificadas'' e apresentadas em destaque, para que a pendência seja corrigida.

\tela[0.8]{indicadores}{Painel de indicadores, página inicial do sistema: posição do período,
evolução de seis meses, inadimplência com os maiores devedores e gastos por categoria.}{indicadores}

\FloatBarrier

% ---------------------------------------------------------------------------
\reqretirado{RF004}{Exportar relatório}{Retirado do escopo do MVP na versão 1.1. O identificador
permanece reservado.}

% ===========================================================================
\section{Administração}

Esta subseção reúne as funções que alteram o comportamento do sistema como um todo. Elas são
restritas à diretoria, por serem decisões de gestão do Centro Acadêmico, e não parte da rotina
financeira diária. Esta é a única subseção do capítulo cujos casos de uso a área administrativa não
executa. As funções ficam reunidas em uma única tela, com uma aba para cada caso de uso, visível
apenas ao perfil de diretoria.

\textbf{Ator de todos os casos de uso desta subseção:} diretoria.

% ---------------------------------------------------------------------------
\req{RF001}{Manter usuários e permissões}

Permite conceder, alterar e revogar o acesso ao sistema.

\prioE

\entradas o usuário autenticado possui perfil de diretoria. Entradas: nome, e-mail, perfil
(administrativo ou diretoria) e situação.

\saidas o usuário é criado, alterado ou inativado; as permissões passam a valer no próximo acesso;
a operação é registrada na trilha de auditoria.

\fluxoprincipal
\begin{enumerate}
  \item O usuário seleciona a opção ``Administração'', aba ``Usuários''
        (Figura~\ref{fig:adm-usuarios}).
  \item O sistema apresenta a relação de usuários, com o perfil, a situação e o último acesso de
        cada um.
  \item O usuário inclui um novo acesso ou seleciona um existente.
  \item O usuário define o perfil, entre administrativo e diretoria.
  \item O sistema grava a alteração e registra a operação.
\end{enumerate}

\fluxossecundarios

\fluxo{A1}{Troca de gestão}
Ao inativar os acessos da diretoria que deixa o cargo, o sistema exige que ao menos um usuário com
perfil de diretoria permaneça ativo, para que o sistema não fique sem administrador.

\fluxo{E1}{Autoinativação}
O sistema recusa que um usuário inative o próprio acesso, evitando o bloqueio acidental do único
administrador ativo.

\tela{administracao-usuarios}{Administração, aba Usuários: perfis, situação e último
acesso.}{adm-usuarios}

\FloatBarrier

% ---------------------------------------------------------------------------
\req{RF002}{Manter categorias de movimentação}

Permite definir as categorias pelas quais entradas e saídas são classificadas, que são o eixo de
análise dos gastos do Centro Acadêmico.

\prioE

\entradas o usuário possui perfil de diretoria. Entradas: nome da categoria, tipo a que se aplica
(entradas, saídas ou ambos) e situação.

\saidas a categoria fica disponível para classificação em [Movimentações.RF001] e passa a compor
os recortes de [Indicadores.RF003].

\fluxoprincipal
\begin{enumerate}
  \item O usuário seleciona a opção ``Administração'', aba ``Categorias''
        (Figura~\ref{fig:adm-categorias}).
  \item O sistema apresenta as categorias existentes, com o número de movimentações associadas a
        cada uma.
  \item O usuário inclui, altera ou inativa uma categoria.
  \item O sistema grava a alteração.
\end{enumerate}

\fluxossecundarios

\fluxo{A1}{Inativação de categoria em uso}
Categoria com movimentações associadas pode ser inativada, deixando de ser oferecida em novos
lançamentos, mas não é excluída. As movimentações anteriores mantêm a classificação, para que os
indicadores de períodos passados continuem reproduzíveis.

\fluxo{E1}{Nome duplicado}
O sistema recusa a criação de categoria com nome igual ao de outra ativa.

\tela{administracao-categorias}{Administração, aba Categorias: tipo, uso e
situação.}{adm-categorias}

\FloatBarrier

% ---------------------------------------------------------------------------
\req{RF003}{Parametrizar cobranças recorrentes}

Permite definir os parâmetros que regem a geração automática das mensalidades e as condições
enviadas ao Asaas, e acompanhar as execuções da geração.

\prioE

\entradas o usuário possui perfil de diretoria. Entradas: valor da mensalidade, dia de vencimento,
dia de geração, percentual de multa por atraso e percentual de juros ao mês.

\saidas os parâmetros passam a valer para as cobranças geradas a partir de então. Cobranças já
emitidas não são afetadas. A alteração fica registrada, com autor, data e hora.

\fluxoprincipal
\begin{enumerate}
  \item O usuário seleciona a opção ``Administração'', aba ``Parâmetros de cobrança''
        (Figura~\ref{fig:adm-parametros}).
  \item O sistema apresenta os parâmetros vigentes, a data e o autor da última alteração e as
        últimas execuções da geração recorrente.
  \item O usuário altera os valores desejados.
  \item O sistema apresenta, para conferência, a data da próxima geração, o vencimento, o número de
        alunos que serão cobrados e o total previsto.
  \item O usuário confirma e o sistema grava os parâmetros.
\end{enumerate}

\fluxossecundarios

\fluxo{A1}{Reajuste da mensalidade}
O novo valor vale apenas para cobranças futuras. O sistema informa explicitamente que as cobranças
já emitidas mantêm o valor anterior, inclusive as vencidas e não pagas.

\fluxo{A2}{Consulta das falhas de uma execução}
A partir de uma execução com falhas, o usuário pode consultar a relação dos alunos que não foram
cobrados e o motivo de cada falha ([Cobranças.RF001], fluxo A2).

\fluxo{E1}{Dia de geração posterior ao vencimento}
O sistema recusa a combinação em que o dia de geração seja posterior ao dia de vencimento, que
produziria cobranças já vencidas no momento da emissão.

\tela{administracao-parametros}{Administração, aba Parâmetros de cobrança: parâmetros vigentes,
conferência da próxima geração e últimas execuções.}{adm-parametros}

\FloatBarrier

% ---------------------------------------------------------------------------
\req{RF004}{Configurar integração com o Asaas}

Permite registrar e atualizar as credenciais que viabilizam a comunicação com a plataforma e
acompanhar a situação da integração.

\prioE

\entradas o usuário possui perfil de diretoria. Entradas: chave de acesso à interface de
programação do Asaas, ambiente de operação (testes ou produção) e prazo para o alerta de
interrupção da conciliação.

\saidas a chave fica armazenada de forma cifrada e passa a ser usada pelos casos de uso de
cobrança; o endereço de notificação do sistema é apresentado para registro na plataforma. A
operação é registrada, sem que a chave seja gravada na trilha.

\fluxoprincipal
\begin{enumerate}
  \item O usuário seleciona a opção ``Administração'', aba ``Integração com o Asaas''
        (Figura~\ref{fig:adm-integracao}).
  \item O sistema apresenta a situação atual da integração (ambiente, último teste, última
        notificação recebida e falhas de entrega seguidas), a data de cadastro da chave e o
        endereço de notificação a ser cadastrado na plataforma.
  \item O usuário aciona ``Substituir chave'' e informa a nova chave e o ambiente
        (Figura~\ref{fig:adm-chave}).
  \item O sistema testa a comunicação com a plataforma.
  \item Confirmada a comunicação, o sistema armazena a chave de forma cifrada e registra a
        operação.
\end{enumerate}

\fluxossecundarios

\fluxo{A1}{Ambiente de testes}
O usuário pode apontar a integração para o ambiente de testes da plataforma, hipótese em que o
sistema sinaliza, em faixa permanente no topo de todas as telas, que as cobranças emitidas não são
reais.

\fluxo{A2}{Prazo do alerta de interrupção}
O usuário escolhe após quantas horas sem notificações a diretoria é alertada (12, 24 ou 48 horas),
respeitado o limite de [Confiabilidade.NF002].

\fluxo{E1}{Chave inválida}
Se o teste de comunicação falhar, o sistema não substitui a chave vigente e informa a falha,
preservando a integração que estava funcionando.

\fluxo{E2}{Exibição da chave}
A chave armazenada não é exibida em tela nem exportada em nenhuma hipótese, sendo apenas
substituível ([Segurança.NF003]).

\tela{administracao-integracao}{Administração, aba Integração com o Asaas: situação da integração,
endereço de notificação e prazo do alerta.}{adm-integracao}
\tela{administracao-dialogo-chave}{Substituição da chave do Asaas: a nova chave é testada antes de
ser salva.}{adm-chave}

\FloatBarrier

% ---------------------------------------------------------------------------
\req{RF005}{Consultar trilha de auditoria}

Permite consultar o registro das operações que alteraram dados no sistema.

\prioD

\entradas o usuário possui perfil de diretoria. Entradas: filtros por período, usuário (incluindo
as rotinas automáticas), tipo de operação e tipo de registro afetado.

\saidas é apresentada a relação de operações registradas, com usuário, data, hora, operação,
registro afetado e detalhe. O sistema não sofre alteração de estado.

\fluxoprincipal
\begin{enumerate}
  \item O usuário seleciona a opção ``Administração'', aba ``Trilha de auditoria''
        (Figura~\ref{fig:adm-trilha}).
  \item O sistema apresenta as operações do período corrente, ordenadas da mais recente para a mais
        antiga.
  \item O usuário ajusta os filtros.
  \item O sistema apresenta o resultado.
\end{enumerate}

\fluxossecundarios

\fluxo{A1}{Consulta a partir de um registro}
A trilha também pode ser aberta a partir de uma cobrança ou movimentação, hipótese em que o filtro
já vem aplicado àquele registro.

\fluxo{E1}{Alteração da trilha}
Não existe operação de alteração ou exclusão de registros da trilha, para nenhum perfil de
usuário. A ausência dessa operação é deliberada e integra a especificação.

\tela{administracao-trilha}{Administração, aba Trilha de auditoria: filtros e operações
registradas, inclusive as das rotinas automáticas.}{adm-trilha}

\FloatBarrier

% ---------------------------------------------------------------------------
% Capítulo 3
% ---------------------------------------------------------------------------

\chapter{Requisitos não funcionais}

Este capítulo apresenta os requisitos não funcionais do IntegraCASD, agrupados por natureza. A
numeração dos identificadores é reiniciada em cada subseção, conforme a convenção descrita na
Introdução.

Duas observações valem para o capítulo inteiro. A primeira é que os requisitos deste capítulo
enunciam propriedades que a solução deve apresentar, e não produtos ou linguagens específicos:
exige-se, por exemplo, que determinadas operações sejam atômicas, e não que se utilize um banco de
dados determinado. A escolha dos meios cabe às decisões de arquitetura, registradas em
\texttt{docs/adr/}. Desde a versão 1.0, foram registradas a adoção de um monólito modular
(ADR-0001) e do PostgreSQL como banco de dados e como fila de tarefas (ADR-0002); onde uma decisão
já atende a um requisito, isso é indicado no próprio requisito.

A segunda é que vários requisitos derivam de características documentadas da plataforma Asaas, da
qual o sistema depende para emitir e conciliar cobranças. Três delas são citadas repetidamente e
convém enunciá-las desde já: as notificações são autenticadas por um token próprio, transmitido no
cabeçalho \texttt{asaas-access-token}; a fila de notificações é interrompida após quinze falhas
consecutivas de entrega, e os eventos não entregues são descartados depois de quatorze dias; e a
conta está sujeita ao limite de vinte e cinco mil requisições a cada doze horas, respondendo com
código 429 ao excedê-lo.

Considera-se, para efeito de dimensionamento, o volume informado pelo CASD de 700 a 800 alunos
cobrados mensalmente.

% ===========================================================================
\section{Usabilidade}

Esta seção descreve os requisitos não funcionais associados à facilidade de uso da interface com o
usuário, material de treinamento e documentação do sistema.

\req{NF001}{Economia de passos nas operações de rotina}
As três operações mais frequentes da área administrativa, que são localizar um aluno, registrar uma
movimentação e emitir uma cobrança avulsa, devem ser alcançáveis em no máximo três interações a
partir da tela inicial do sistema, sem exigir navegação por menus intermediários. No protótipo, o
requisito é atendido pela barra superior, presente em todas as telas, com a busca de aluno e os
botões ``Nova cobrança'' e ``Nova movimentação''.
\prioI
\casos{[Cadastros.RF001], [Cobranças.RF002], [Movimentações.RF001].}

\req{NF002}{Mensagens de erro orientadas à ação}
Toda mensagem de erro apresentada ao usuário deve informar o que deve ser feito para resolver a
situação, e não apenas o que falhou. É vedada a exibição de código de erro, de mensagem originada
da interface de programação do Asaas ou de qualquer termo técnico de integração nas telas de uso
corrente, cabendo ao sistema traduzir essas ocorrências para linguagem compreensível pela área
administrativa.
\prioI

\req{NF003}{Confirmação explícita em operação irreversível}
As operações que não podem ser desfeitas devem exigir etapa adicional de confirmação, na qual o
sistema apresente um resumo do efeito da operação antes de executá-la. Aplica-se ao cancelamento de
cobrança, ao registro manual de pagamento, ao estorno de movimentação e à inativação de aluno com
débito.
\prioE
\casos{[Cadastros.RF001], [Cobranças.RF004], [Cobranças.RF006], [Movimentações.RF002].}

\req{NF004}{Informação que não depende apenas de cor}
Nenhuma informação pode ser transmitida apenas por cor. As situações de cadastros, cobranças e
movimentações são indicadas por selos com texto, acompanhados de ícone quando ajuda; os gráficos
usam paleta legível por pessoas com daltonismo, apresentam o valor ao passar o mouse e oferecem os
mesmos dados em tabela.
\prioI
\casos{[Indicadores.RF001], [Indicadores.RF003].}

% ===========================================================================
\section{Confiabilidade}

Esta seção descreve os requisitos não funcionais associados à frequência e severidade de falhas do
sistema e à habilidade de recuperação das mesmas, bem como à correção do sistema.

\req{NF001}{Disponibilidade do endereço de notificação}
O endereço que recebe as notificações do Asaas deve apresentar disponibilidade mínima de 99\% ao
mês e, em nenhuma hipótese, acumular quinze falhas consecutivas de resposta, que é o limite a
partir do qual a plataforma interrompe o envio da fila de eventos. A exigência aplica-se inclusive
nos dias de vencimento das mensalidades, quando se concentra o maior volume de notificações.
\prioE
\casos{[Cobranças.RF005].}

\req{NF002}{Alerta de interrupção da conciliação}
O sistema deve emitir alerta à diretoria financeira do CASD quando deixar de receber notificações
do Asaas por período superior ao configurado ([Administração.RF004]), não excedendo quarenta e oito
horas. O alerta destina-se a permitir a correção da falha dentro do prazo de quatorze dias em que a
plataforma retém os eventos não entregues, findo o qual eles são descartados em definitivo e os
pagamentos correspondentes deixam de ser conciliáveis.
\prioE
\casos{[Cobranças.RF005], [Cobranças.RF007], [Administração.RF004].}

\req{NF003}{Idempotência no processamento de notificações}
O processamento de uma notificação já recebida não pode produzir efeito adicional. O sistema deve
identificar a repetição pelo identificador do evento no Asaas e descartá-la, de modo que nenhuma
cobrança seja alterada duas vezes e nenhuma movimentação de entrada seja duplicada, ainda que a
mesma notificação seja entregue diversas vezes pela plataforma.
\prioE
\casos{[Cobranças.RF005].}

\req{NF004}{Atomicidade da baixa de cobrança}
A atualização da situação de uma cobrança e a criação da movimentação de entrada correspondente
devem ocorrer de forma atômica, isto é, ambas as operações se completam ou nenhuma delas produz
efeito. Não é admissível que o sistema registre uma cobrança como paga sem a movimentação
correspondente, nem o contrário, situação que faria o painel financeiro divergir do histórico de
cobranças.
\prioE
\casos{[Cobranças.RF005], [Cobranças.RF006], [Cobranças.RF007].}

\req{NF005}{Correspondência entre o sistema e a plataforma}
Nenhuma cobrança pode existir apenas no IntegraCASD. Toda cobrança registrada no sistema corresponde
a uma cobrança criada no Asaas, e a falha na criação junto à plataforma impede o registro local,
conforme os fluxos de exceção especificados no Capítulo~2.
\prioE
\casos{[Cobranças.RF001], [Cobranças.RF002].}

\req{NF006}{Cópia de segurança e recuperação}
O sistema deve dispor de cópia de segurança diária da base de dados, retida por no mínimo trinta
dias. O procedimento de recuperação deve ser executado e documentado ao menos uma vez antes da
entrega do sistema, de modo que a sua eficácia não seja apenas presumida.
\prioE

\req{NF007}{Registro das execuções automáticas}
Cada execução das rotinas que operam sem intervenção humana deve gravar data e hora de início e de
término, o total de registros processados e a relação dos que falharam, de modo que a área
administrativa possa verificar se a rotina foi executada e o que dela resultou.
\prioI
\casos{[Cobranças.RF001], [Cobranças.RF007], [Administração.RF003].}

% ===========================================================================
\section{Desempenho}

Esta seção descreve os requisitos não funcionais associados à eficiência, uso de recursos e tempo
de resposta do sistema.

\req{NF001}{Tempo de resposta ao Asaas}
O sistema deve responder a toda notificação recebida em até cinco segundos, confirmando o
recebimento antes de processar o seu conteúdo. O processamento deve ocorrer em seguida, de forma
desacoplada da resposta, de modo que o tempo de tratamento interno não influencie o tempo de
resposta à plataforma. Este requisito decorre diretamente de [Confiabilidade.NF001]: resposta lenta
é contabilizada como falha de entrega pelo Asaas. A solução adotada (gravar o evento, enfileirá-lo
e responder) está descrita no ADR-0002.
\prioE
\casos{[Cobranças.RF005].}

\req{NF002}{Tempo de resposta das operações interativas}
Em condições normais de operação, 95\% das operações interativas devem ser respondidas em até três
segundos, e nenhuma deve ultrapassar dez segundos. A apuração de indicadores está dispensada deste
limite, sujeitando-se ao requisito próprio abaixo, assim como as operações que dependem de resposta
do Asaas.
\prioI

\req{NF003}{Tempo de apuração dos indicadores}
O painel de indicadores deve ser apresentado em até cinco segundos para um período de referência de
até doze meses, considerado o volume de movimentações decorrente de 800 alunos cobrados
mensalmente.
\prioI
\casos{[Indicadores.RF001], [Indicadores.RF002], [Indicadores.RF003].}

\req{NF004}{Janela de execução da geração recorrente}
A geração das cobranças recorrentes deve concluir em até trinta minutos para uma base de até mil
alunos ativos, consumindo no máximo dez por cento da cota de vinte e cinco mil requisições a cada
doze horas atribuída à conta do CASD. O sistema deve espaçar as requisições à plataforma de modo a
não exceder os limites de frequência e de concorrência da interface de programação.
\prioI
\casos{[Cobranças.RF001].}

\req{NF005}{Tratamento de excesso de requisições}
Ao receber da plataforma resposta com código 429, indicativa de excesso de requisições, o sistema
deve aguardar o intervalo informado pela própria resposta e repetir a requisição, sem perder a
operação em curso nem interromper a rotina que a originou.
\prioI
\casos{[Cobranças.RF001], [Cobranças.RF007].}

% ===========================================================================
\section{Segurança}

Esta seção descreve os requisitos não funcionais associados à integridade, privacidade e
autenticidade dos dados do sistema.

\req{NF001}{Autenticação e perfis concêntricos}
Todo acesso ao sistema deve ser autenticado. As permissões são concedidas por perfil, sendo que o
perfil de diretoria possui todas as permissões do perfil administrativo, acrescidas das funções
reunidas na subseção Administração do Capítulo~2. Nenhuma função é exclusiva do perfil
administrativo.
\prioE
\casos{[Administração.RF001].}

\req{NF002}{Autenticação das notificações recebidas}
Toda notificação recebida da plataforma deve ter o cabeçalho \texttt{asaas-access-token}
verificado antes de qualquer processamento. Notificação sem token válido é rejeitada, não produz
efeito algum sobre os dados e é registrada como tentativa de acesso indevido. O token utilizado
nessa verificação deve ser distinto da chave de acesso à interface de programação, conforme orienta
a documentação da plataforma.

Este é o único ponto do sistema exposto na Internet capaz de alterar dados financeiros sem usuário
autenticado, o que confere ao requisito importância superior à que o seu enunciado sugere.
\prioE
\casos{[Cobranças.RF005].}

\req{NF003}{Guarda da chave de acesso à plataforma}
A chave de acesso à interface de programação do Asaas deve ser armazenada de forma cifrada, não
pode ser exibida em tela, exportada, transmitida por qualquer meio ao usuário nem gravada em
registro de execução, em log ou na trilha de auditoria. A sua substituição é a única operação
admitida sobre ela. Os demais segredos do sistema ficam apenas na configuração do ambiente, nunca no
código-fonte.
\prioE
\casos{[Administração.RF004].}

\req{NF004}{Proteção da comunicação}
Toda comunicação com o sistema deve ocorrer sobre HTTPS, com TLS na versão 1.2 ou superior,
inclusive no endereço que recebe as notificações da plataforma.
\prioE

\req{NF005}{Proteção das credenciais dos usuários}
As senhas dos usuários devem ser armazenadas exclusivamente na forma de resumo criptográfico com
sal, calculado por algoritmo próprio para esse fim, sendo vedado o armazenamento em texto claro ou
por meio reversível.
\prioE

\req{NF006}{Trilha de auditoria inalterável}
Toda operação que altere cadastro, cobrança, movimentação ou configuração deve ser registrada com a
identificação do usuário (ou da rotina automática), data, hora, operação realizada e registro
afetado. Não deve existir, para nenhum perfil de usuário, função de alteração ou de exclusão dos
registros da trilha.
\prioI
\casos{[Administração.RF005], [Cobranças.RF006], [Movimentações.RF002].}

\req{NF007}{Bloqueio após tentativas de acesso malsucedidas}
Cinco tentativas consecutivas de autenticação malsucedidas devem bloquear temporariamente o acesso
àquela identificação de usuário, registrando-se a ocorrência.
\prioI

% ===========================================================================
\section{Distribuição}

Esta seção descreve os requisitos não funcionais associados à distribuição da versão executável do
sistema.

\req{NF001}{Distribuição exclusivamente web}
O sistema deve ser distribuído como aplicação web, sem exigir instalação de qualquer componente nas
estações de trabalho do CASD ou nos dispositivos dos usuários. As atualizações devem ser aplicadas
no ambiente servidor e refletidas para todos os usuários sem ação da parte deles.
\prioE

\req{NF002}{Carga inicial a partir das planilhas}
O sistema deve ser implantado com a base de dados já carregada a partir das planilhas hoje
utilizadas pelo CASD. A carga é procedimento de implantação, executado uma única vez por script
próprio, e não funcionalidade do sistema, razão pela qual não há tela destinada a ela. Como as
planilhas trazem dados fora do padrão (CPF sem zeros à esquerda, valores em texto, situações em
texto livre), a carga deve normalizá-los e relacionar à parte os registros que não puder
interpretar. A área administrativa deve conferir e aprovar formalmente a carga antes da entrada em
operação, uma vez que os dados eventualmente omitidos terão de ser digitados manualmente.
\prioE

\req{NF003}{Retorno à versão anterior}
O procedimento de implantação deve permitir o retorno à versão imediatamente anterior em até
trinta minutos, preservando os dados registrados até o momento da falha.
\prioD

% ===========================================================================
\section{Padrões}

Esta seção descreve os requisitos não funcionais associados a padrões ou normas que devem ser
seguidos pelo sistema ou pelo seu processo de desenvolvimento.

\req{NF001}{Conformidade com a Lei Geral de Proteção de Dados}
O tratamento dos dados pessoais mantidos pelo sistema deve observar a Lei nº 13.709/2018, figurando
o Centro Acadêmico Santos Dumont como controlador dos dados e o Asaas como operador, na parte que
lhe cabe. Aplicam-se, em particular, as seguintes regras:
\begin{itemize}
  \item a coleta limita-se ao necessário para identificar o aluno e para emitir e conciliar as
        cobranças; por essa razão, a data de nascimento não é coletada;
  \item dados reais de alunos nunca entram no repositório do projeto, nos testes ou nos exemplos,
        que usam apenas dados fictícios, e o ambiente de homologação usa dados anonimizados
        ([Hardware e software.NF004]);
  \item o sistema deve permitir atender à solicitação de acesso formulada pelo titular e, quanto à
        solicitação de eliminação, deve conciliá-la com a regra de que nada é excluído fisicamente:
        os dados pessoais do aluno inativo são anonimizados, enquanto os registros financeiros são
        conservados, nas hipóteses de guarda previstas no art.~16 da Lei.
\end{itemize}
\prioE
\casos{[Cadastros.RF001].}

\req{NF002}{Aderência à documentação da plataforma}
A integração com o Asaas deve seguir a documentação oficial da plataforma quanto à autenticação, ao
tratamento de erros e aos eventos de cobrança, sendo vedado depender de comportamento não
documentado ou observado apenas empiricamente.
\prioI

\req{NF003}{Padrão de codificação e integração}
A equipe deve adotar um padrão de codificação único, registrado no repositório do projeto, e só
integrar uma alteração ao ramo principal com a verificação estática, os testes e a compilação
passando. O padrão está registrado em \texttt{CONTRIBUTING.md}: ramos a partir do principal,
mensagens de \textit{commit} no formato Conventional Commits e integração por \textit{merge} no ramo
principal.
\prioD

% ===========================================================================
\section{Hardware e software}

Esta seção descreve os requisitos não funcionais associados ao hardware e software usados para
desenvolver ou para executar o sistema.

\req{NF001}{Endereço público e estável para a conciliação}
O ambiente de produção deve dispor de endereço público, estável e alcançável pela Internet,
protegido por HTTPS, destinado a receber as notificações da plataforma. A hospedagem dar-se-á em
serviço de nuvem contratado ou em infraestrutura fornecida pelo ITA, conforme decisão a ser
registrada no ADR-0006.

Trata-se de pré-requisito de infraestrutura sem o qual a conciliação automática não opera, e cuja
indisponibilidade não é contornável por programação.
\prioE
\casos{[Cobranças.RF005].}

\req{NF002}{Empacotamento em contêineres}
O sistema deve ser empacotado em contêineres, tanto no ambiente de desenvolvimento quanto nos de
homologação e produção, de modo a manter a equivalência entre eles e a reduzir a dependência da
configuração das máquinas dos integrantes da equipe.
\prioI

\req{NF003}{Uso do ambiente de testes da plataforma}
Todo o desenvolvimento e toda a homologação devem utilizar o ambiente de testes do Asaas, jamais o
ambiente de produção. Quando o sistema estiver apontado para o ambiente de testes, deve exibir
sinalização visível e permanente de que as cobranças emitidas não são reais. O ambiente de
desenvolvimento deve, ainda, ser alcançável pelas notificações da plataforma, o que exige endereço
público temporário ou equivalente, já que contêineres executados em máquina local não são
acessíveis pela Internet.
\prioE
\casos{[Cobranças.RF005], [Administração.RF004].}

\req{NF004}{Ambiente de homologação}
Deve ser mantido ambiente de homologação equivalente ao de produção, com dados anonimizados,
destinado à validação das entregas pela diretoria do CASD antes de cada implantação.
\prioI

\req{NF005}{Propriedades exigidas do armazenamento}
O mecanismo de armazenamento adotado deve suportar transações que abranjam mais de um registro, de
modo a viabilizar o disposto em [Confiabilidade.NF004], e permitir a realização de cópias de
segurança consistentes sem interrupção do serviço. O ADR-0002 adota o PostgreSQL, que atende a
essas propriedades.
\prioE
\casos{[Cobranças.RF005].}

% ---------------------------------------------------------------------------
% Capítulo 4
% ---------------------------------------------------------------------------

\chapter{Descrição da interface com o usuário}

Este capítulo descreve a organização das telas do IntegraCASD e as decisões de interface que valem
para todas elas. As telas foram desenhadas em um protótipo navegável, com dados fictícios e sem
comunicação com o servidor, disponível em \texttt{docs/telas/}. As imagens de cada tela estão no
Capítulo~2, junto aos casos de uso que ilustram.

\section{Organização das telas}

Todas as telas compartilham a mesma estrutura, visível na Figura~\ref{fig:cobranca}:
\begin{itemize}
  \item \textbf{Menu lateral}, com uma entrada para cada subseção do Capítulo~2: Indicadores,
        Cadastros, Cobranças, Movimentações e Administração. A entrada de Administração traz a
        indicação ``Diretoria'' e só aparece para esse perfil. No rodapé do menu ficam o nome e o
        cargo do usuário autenticado.
  \item \textbf{Barra superior}, com a busca de aluno por nome ou CPF e os botões ``Nova
        movimentação'' e ``Nova cobrança'', que abrem o formulário em uma janela sobre a tela atual,
        sem perder o contexto ([Usabilidade.NF001]).
  \item \textbf{Faixa de ambiente de testes}, no topo de todas as telas, enquanto a integração
        estiver apontada para o ambiente de testes do Asaas ([Hardware e software.NF003]).
\end{itemize}

A página inicial é o painel de indicadores. Cada tela corresponde a uma pasta de funcionalidade no
código do frontend, e as telas com mais de um tipo de registro (Cadastros e Administração) usam
abas.

\section{Decisões de interface}

\begin{itemize}
  \item \textbf{Confirmação com resumo do efeito} no cancelamento de cobrança, no registro manual
        de pagamento, no estorno de movimentação e na inativação de aluno com débito
        ([Usabilidade.NF003]).
  \item \textbf{Selos sempre com texto} e, quando ajuda, ícone; nunca apenas cor
        ([Usabilidade.NF004]).
  \item \textbf{Amarelo do CASD apenas} no botão principal de cada tela, no item ativo do menu e em
        alertas, para que chame a atenção onde ela é necessária.
  \item \textbf{Gráficos} com paleta validada para daltonismo (azul para entradas, laranja para
        saídas), valor ao passar o mouse e os mesmos dados também em tabela.
  \item \textbf{Valores sempre originais} nas cobranças, com aviso de que multa e juros são
        aplicados pelo Asaas.
  \item \textbf{O detalhe da cobrança é o centro das ações}; a lista de cobranças apenas mostra e
        leva ao detalhe.
  \item \textbf{O estorno é feito direto da lista} de movimentações, sem tela de detalhe; nas
        movimentações originadas de cobrança, o botão fica desabilitado e uma dica explica onde
        corrigi-las.
  \item \textbf{A lista de cobranças abre com as abertas e vencidas}, ordenadas por vencimento; as
        demais situações são acessadas pelos cartões de resumo.
\end{itemize}

A tela de autenticação não foi desenhada no protótipo, por seguir o padrão usual de identificação
e senha; os requisitos que a regem são [Segurança.NF001], [Segurança.NF005] e [Segurança.NF007].

\section{Correspondência entre telas e casos de uso}

\begin{longtable}{@{}p{0.2\linewidth} p{0.46\linewidth} >{\raggedright\arraybackslash}p{0.26\linewidth}@{}}
  \toprule
  \textbf{Tela} & \textbf{O que mostra} & \textbf{Requisitos e figuras} \\
  \midrule
  \endhead
  Indicadores & Página inicial: entradas, saídas, resultado e valor aguardando liberação; evolução
                de seis meses; inadimplência com atalho para as cobranças vencidas; gastos por
                categoria &
                Indicadores.RF001 a RF003; Fig.~\ref{fig:indicadores} \\
  Cadastros & Abas de alunos, funcionários, contratos e projetos; formulário de novo aluno &
              Cadastros.RF001, RF002, RF004 e RF005; Figs.~\ref{fig:cadastros-alunos}
              a~\ref{fig:cadastros-novo-aluno}, \ref{fig:cadastros-funcionarios},
              \ref{fig:cadastros-contratos} e~\ref{fig:cadastros-projetos} \\
  Aluno & Dados editáveis, cobranças, movimentações, nova cobrança com o aluno preenchido e
          inativação com aviso de débito & Cadastros.RF001; Figs.~\ref{fig:aluno-dados}
          a~\ref{fig:aluno-inativar} \\
  Cobranças & Estado da conciliação, resumo por situação, filtros, lista de abertas e vencidas e
              sincronização geral & Cobranças.RF007 e RF008; Figs.~\ref{fig:cobrancas}
              e~\ref{fig:cobrancas-sincronizar} \\
  Cobrança & Detalhe: boleto, sincronização, pagamento manual, cancelamento e
             renegociação, histórico de eventos do Asaas & Cobranças.RF003 a RF007;
             Figs.~\ref{fig:cobranca}, \ref{fig:cobranca-cancelar},
             \ref{fig:cobranca-pagamento-manual} e~\ref{fig:cobranca-sincronizar} \\
  Movimentações & Totais do recorte, filtros, estorno e movimentação estornada com o seu espelho &
                  Movimentações.RF002 e RF003; Figs.~\ref{fig:movimentacoes-estornar}
                  e~\ref{fig:movimentacoes} \\
  Barra superior & Busca de aluno, nova cobrança avulsa e nova movimentação & Cobranças.RF002,
                   Movimentações.RF001; Figs.~\ref{fig:nova-cobranca}
                   e~\ref{fig:nova-movimentacao} \\
  Administração & Usuários, categorias, parâmetros com conferência e execuções, integração com o
                  Asaas e trilha de auditoria & Administração.RF001 a RF005;
                  Figs.~\ref{fig:adm-usuarios} a~\ref{fig:adm-trilha} \\
  \bottomrule
\end{longtable}

% ---------------------------------------------------------------------------
% Apêndice
% ---------------------------------------------------------------------------

\appendix
\chapter{Pontos em aberto}

As questões abaixo ainda dependem de decisão da equipe ou do CASD e serão resolvidas nas próximas
versões deste documento.

\begin{enumerate}
  \item \textbf{Confirmação do recebimento no Asaas.} No registro manual de pagamento
        ([Cobranças.RF006]), avaliar se o sistema deve também informar o recebimento ao Asaas, para
        que o boleto perca a validade e as duas bases não divirjam. Na especificação atual, a
        divergência é aceita e fica sinalizada pela situação ``Paga com registro manual''.
  \item \textbf{Apresentação das pendências de conciliação.} Definir se as notificações sem
        cobrança correspondente, as divergências da sincronização e as falhas da geração recorrente
        serão reunidas em uma tela própria de pendências ou se continuam apresentadas apenas no
        resultado da sincronização e no histórico de execuções.
  \item \textbf{Alunos bolsistas ou isentos.} Definir como são tratados na geração recorrente.
  \item \textbf{Situações acadêmicas da planilha.} Definir quais situações (por exemplo, trancado
        ou em intercâmbio) correspondem a aluno ativo, isto é, cobrado na geração recorrente.
  \item \textbf{Histórico de boletos anteriores.} Definir se as cobranças emitidas antes da
        implantação entram na carga inicial.
  \item \textbf{Resultado da busca na barra superior.} Definir se a busca abre uma lista de
        resultados ou vai direto ao cadastro do aluno quando houver um único resultado.
\end{enumerate}

\end{document}
